\documentclass[11pt]{article}
\usepackage[letterpaper,margin=1in]{geometry}
\usepackage{amsmath,enumitem,xcolor,booktabs,longtable,array,verbatim}
\usepackage[hidelinks]{hyperref}
\setlength{\parskip}{0.35em}\setlength{\parindent}{0pt}
\setlist{itemsep=0.1em,topsep=0.2em,leftmargin=1.6em}
\setlength{\emergencystretch}{3em}
\definecolor{place}{RGB}{176,58,46}
\definecolor{open}{RGB}{31,97,141}
\definecolor{have}{RGB}{30,110,60}
\definecolor{miss}{RGB}{176,58,46}
\definecolor{stale}{RGB}{160,100,0}
\definecolor{cut}{RGB}{110,60,150}
\newcommand{\Place}[2]{\par\smallskip\noindent\fbox{\parbox{0.97\linewidth}{%
  \textcolor{place}{\textbf{#1}}\quad #2}}\par\smallskip}
\newcommand{\Open}[1]{\textcolor{open}{\textit{Open: #1}}}
\newcommand{\Have}{\textcolor{have}{\textbf{[written]}}\ }
\newcommand{\Miss}{\textcolor{miss}{\textbf{[to write]}}\ }
\newcommand{\Stale}{\textcolor{stale}{\textbf{[fix]}}\ }
\newcommand{\Cut}{\textcolor{cut}{\textbf{[trim]}}\ }
\newcommand{\Note}{\textcolor{open}{\textbf{[from the notes]}}\ }
\newcommand{\Cite}[1]{\par{\footnotesize\raggedright\textcolor{have}{\textbf{Cite:}} #1\par}}
% a paragraph of the draft: title, where it is, one-line overview
\newcommand{\Pg}[3]{\item \textbf{#1} {\small(#2)}\par\textit{#3}}
% kinds of suggested change
\newcommand{\tC}{\textcolor{miss}{\textbf{[correct]}}\ }
\newcommand{\tG}{\textcolor{stale}{\textbf{[wording]}}\ }
\newcommand{\tT}{\textcolor{cut}{\textbf{[trim]}}\ }
\newcommand{\tA}{\textcolor{have}{\textbf{[add]}}\ }
\newcommand{\Msg}[1]{\par\textcolor{open}{\textbf{Message:}} \textit{#1}\par}
\newcommand{\Para}[2]{\textbf{#1} \textnormal{\small(L#2)}}
\newcommand{\New}[1]{\textbf{#1} \textnormal{\small(new)}}
% ---- literature check (2026-10-15): marks, snippets, draft macros shown as text
\definecolor{chk}{RGB}{0,90,170}
\newcommand{\add}[1]{\textcolor{chk}{#1}}
\newcommand{\del}[1]{\textcolor{red}{[del: #1]}}
\newcommand{\Chk}[1]{\textcolor{chk}{\textbf{[lit.\ check #1]}}\ }
\newcommand{\Snippet}[1]{\par\smallskip\noindent\colorbox{chk!8}{\parbox{0.96\linewidth}{\small\textit{Suggested text:} #1}}\par\smallskip}
\providecommand{\citep}[1]{(\texttt{\detokenize{#1}})}
\providecommand{\citet}[1]{\texttt{\detokenize{#1}}}
\providecommand{\cref}[1]{\texttt{\detokenize{#1}}}
\providecommand{\alphaC}{\alpha_c}
\sloppy
\title{Manuscript outline: what each paragraph should contain}
\author{enceladus\_DSM \quad \texttt{studies/keff\_sintering/manuscript\_outline\_new.tex}}
\date{Re-synced 2026-10-15 against the Overleaf project uploaded that day (second upload), with the literature check merged in}
\begin{document}
\maketitle

% ===================== Literature check, 2026-10-15 ==========================
% Merged from three independent reviews (Literature/review/{A,B,C}); every
% finding below was re-checked by the coordinator against the cited page.
% Draft locations are .tex source lines of the 10-15 draft (``tex L''); it
% compiles with finalold, so it has no printed line numbers.
\section*{\add{Literature check of the 10-15 draft: ranked index}}
\add{Three independent reviews (\texttt{Literature/review/\{A,B,C\}}), merged
and re-checked page by page against the library. The evidence (paper, PDF
page, short quote) and any \emph{suggested text} sit in the body, at the
entries marked \Chk{F$n$}. \textbf{F1}--\textbf{F5} change what the paper
claims or how it justifies a choice; settle them first.}
\begin{enumerate}[label=\textbf{F\arabic*},leftmargin=2.2em]
\item \add{\textbf{The gap must be stated against Bouvet et al.\ (2022)}, who
  computed $k_{\mathrm{eff}}$ through isothermal metamorphism; the gradient
  sentence is Calonne et al.\ (2014), not 2011. (Intro \P8.)}
\item \add{\textbf{$\alpha_c = 10^{-3}$ is attachment-limited, not
  diffusion-limited.} (Intro \P6; Section~3.2.)}
\item \add{\textbf{Surface diffusion does not explain the $-20\,^\circ$C
  shortfall}, and the model gives 41\% \emph{of} the growth. (Section~5.1.)}
\item \add{\textbf{Exponent sources}: Hobbs \& Mason give $1/5$ for vapor
  through air; Thomas et al.\ give 4--5 for surface diffusion; the exponent
  alone does not identify the mechanism; Kingery's $1/7$ is disputed;
  ``planetary timescales'' reverses Molaro. (Intro \P5; Section~5.1.)}
\item \add{\textbf{The 2~mm domain cannot rest on Calonne et al.\ (2011)},
  whose REVs are 2.5--5.5~mm; use our domain-size test. (Section~2.5,
  Aggregates.)}
\item \add{\textbf{New \P4}: ``at least 180~K''; ``strengthening''; the
  0.01--1~mm\,yr$^{-1}$ range is Choukroun's summary; burial and
  conductivity sentences need citations. (Intro \P4.)}
\item \add{\textbf{Bouvet et al.\ (2023) $\to$ (2022)} for the time
  rescaling through $\alpha_c$. (Section~5.1, $\alpha_c$ as an effective
  parameter.)}
\item \add{\textbf{$\alpha_c$ range and form}: cite Kaempfer \& Plapp and
  Braun et al.; not Arrhenius. (Intro \P6.)}
\item \add{\textbf{Equations}: $P_a$ vs $P_{\mathrm{atm}}$; a stray $T_0$;
  $k_j$ source; ``ASHRAE''; empty scaling appendix (\texttt{Sclaing});
  Eq.~(11d) missing $\nabla$. (Sections~2.1, 2.3.)}
\item \add{\textbf{Smaller citation fixes}: Colbeck (1998) for depth hoar;
  Postberg's 6\%; Swinkels \& Ashby; Colbeck (1983) for experiments; Flin
  for bulk properties; Auriault et al.\ (2009) book; Nicoli's scalar law.
  (Intro \P2, \P3, \P5, \P7, \P8; Section~2.)}
\item \add{\textbf{Bibliography}: Byrne et al.\ (2024) is the preprint (use
  doi:10.1029/2024JE008615); \texttt{noauthor\_sintering\_nodate} duplicates
  \texttt{Hobbs\_and\_Mason}; Molaro pages 243--276; Calonne (2011) L23501;
  Wexler pages 5--20; Flin ``2014'' is \emph{Physics and Chemistry of Ice
  2010}, Hokkaido Univ.\ Press, 2011, pp.~321--328.}
\item \add{\textbf{Typos} in the new text are listed under each paragraph.}
\end{enumerate}
\add{\textbf{Citations worth adding} (closely matching, checked):
Legagneux \& Domin\'e (2005, p.9), whose vapor-transport model shows
temperature as ``a translation of the curve along the time axis'' (for
Section~5.2); Calonne et al.\ (2019, p.5), bonds ``act as heat paths'' at low
density (for Section~5.3). \textbf{Not checkable here:} Hobbs \& Mason
(1964), Massman (1998), Moure \& Fu's Supporting Information.}
% ===================== end literature check ==================================


\textbf{Framing.} Dry snow metamorphism is the story; the model covers one
part of it, sintering by vapor transport, and the bulk conductivity that
results, under conditions a laboratory can reach. Enceladus is the motivation
and the closing implication. Every claim stays inside the simulated range:
$-40$ to $-5\,^{\circ}$C, air at 1~atm, sealed pores, 50~$\mu$m grains,
30~days, two dimensions.

\textbf{How to read this.} Every paragraph of the draft has a numbered
entry: its subject, where it is, a one-line overview, and bullets for what I
suggest changing. Line numbers (L\,first--last) are those of the source of
2026-10-14 compiled here (31 pages; the body ends at L492). That upload
changed only the Introduction, so the Introduction entries are new from a
second full read, and every later entry is the 10-13 review with its line
numbers moved up by 16. Each bullet is tagged:
\tC a statement that is wrong or unsupported;
\tG grammar or wording;
\tT text that can go or shrink;
\tA something missing.
A paragraph with no \tC bullet says nothing I found to be wrong. In the
Introduction, \tA bullets marked \Note or ``(discussed 10-xx)'' are points we
have already worked out or that the old Discussion carried, and which the
draft does not yet contain. Figure captions are reviewed in their own
section, after Section~6. Entries
marked ``to write'' do not exist yet; their bullets are the content. Boxes
mark where a figure or table goes. From Section~5.2, paragraph~5, on, the
entries keep the earlier form (a ``Message'' and a content list), since
nothing is drafted there.

\textbf{First round (user, 2026-10-14).} This round checks that what the
draft says is correct, free of typos, not missing a major argument, well
ordered and concise. A paragraph that does all that is marked ``no changes'';
fine polish waits for the second round, after the team has read it.

\textbf{The old Discussion is gone from the source.} Its literature points
are kept here, marked \Note. Their line numbers refer to the PDF of
2026-10-12 in \texttt{Literature/}, which still contains that text. Each
point is to be checked against its source before use.

\textbf{Suggested citations} follow each paragraph, with the key in your
\texttt{.bib}. Most are the ones the old Discussion used. I have matched them
by title only; check each claim against the paper before it goes in.

\section*{Decisions already made, which this plan follows}

Collected from the work of 2026-09-08 to 2026-10-12: the session log, the
campaign record and the manuscript plan.
\begin{itemize}
\item \textbf{Framing} (10-08, 10-10). Laboratory-observable sintering by
  vapor transport, inside the dry-snow-metamorphism story. The Introduction
  and title stay as written; scope caveats go in the Methods. No timescales
  for Enceladus temperatures; the timescale map is dropped. Advisor's
  comments still to come.
\item \textbf{Comparisons} (10-08). Not natural snow (angular, more porous,
  loaded). The analogues are unloaded ice-sphere and ice-powder samples:
  Molaro et al., Demmenie et al., Choukroun et al. From the snow literature
  only functional forms carry over, never magnitudes.
\item \textbf{Kinetics} (09-22, 10-09). $\alpha_c$ constant in temperature;
  the thin-interface term always on; aggregates stated as $\alpha_c = 1.020
  \times 10^{-3}$. No further simulations for this paper.
\item \textbf{Two-grain argument} (09-09, 10-02, 10-13). Rests on what the
  vapor route does and does not reproduce, and on exponents; no Thomas or
  Kingery runs. The two mismatches have different causes and neither is
  attributed to missing surface diffusion (Appendix, this outline): the rate
  at $-20\,^{\circ}$C is a matter of how fast vapor can be supplied, in an
  experiment whose driving force is a few millikelvin-equivalent; the
  exponent at saturation now matches: from grains just touching the model
  gives 0.29, inside Demmenie's 0.26--0.33 and on the sharp-interface law
  (10-15); the earlier 0.21 came from a pre-made starting neck. Molaro's conventions: neck \emph{width}, time anchored on the first
  measurement. The $-20\,^{\circ}$C series is the benchmark.
\item \textbf{Conductivity} (09-23). Tensor law in the interface; the cell
  problem is a steady problem on a frozen microstructure and says nothing
  about the transient model. Two dimensions only, so no absolute
  $k_{\mathrm{eff}}$ is claimed.
\item \textbf{Clock and normalization} (10-08, 10-12). The clock is
  $\tau_{\mathrm{sub}}$; curves are normalized at $\theta_r = 30$. The first day depends on the interface parameter, so a
  ``rise from day 1'' is never quoted as a property of the ice.
\item \textbf{Notation} (09-28 to 09-30). $\phi$ phase field, $\varphi$
  porosity, $\zeta$ indicator, $\chi$ curvature; subscripts ``ice'' and
  ``air'', not $i$ and $a$; no index notation; a single dot for every
  contraction; the vector $\mathbf{t}$ named as in Auriault et al.
\item \textbf{Citations} (09-23). Only work whose setting matches ours;
  every DOI checked against Crossref.
\item \textbf{Words.} Desublimation (``deposition'' is plume fallout);
  ``vapor''; $\varepsilon$ is the interface parameter, not the interface
  width.
\end{itemize}

\section*{Where the draft stands}

\begin{center}\begin{tabular}{p{0.34\linewidth}p{0.58\linewidth}}\toprule
Section & State\\\midrule
Key points, Abstract, Summary & to write (placeholders)\\
1 Introduction & third draft (10-15): new \P4 on Enceladus conductivity; open: the gap against Bouvet et al.\ (2022), the exponent sources, $\alpha_c$ regime, why SSA\\
2.1--2.3 Model, kinetics, homogenization & written; \textbf{$T$ back to $^\circ$C (10-16)}, $T + T_0$ where kelvin is needed; scaling appendix empty; Eqs.~(2), (11d) to fix\\
2.4, 2.5 Numerics, simulation design & headings only\\
3 Grain pairs & written; 3.2 to rewrite for the tangent-start run ($a = 0.29$, Figure~4 rebuilt 10-15)\\
4 Aggregates & written; numbers to add, one paragraph to correct\\
5.1 What vapor transport reproduces & three paragraphs drafted; the first two need revising for what was found on 10-13\\
5.2 Temperature as a rescaling of time & one paragraph begun\\
5.3--5.6, 6 Conclusion & headings only\\
Appendices A, B & headings only; content is in the Appendix of this outline\\\bottomrule
\end{tabular}\end{center}

\textbf{Done since the 10-14 upload} (the 10-15 drafts): new Introduction
\P4 (Enceladus deposits and conductivity); \P5 rebuilt around the exponents
with Hobbs \& Mason, Kingery, Thomas and Demmenie in the \texttt{.bib}; \P6
with Libbrecht \& Rickerby, layer nucleation and a constant $\alpha_c$; \P7
with Braun et al.\ described; \P8 ending on the paper's gap; \P9's two findings
restated; Figure~1 caption vapor-only; kelvin throughout Section~2.1 (reversed to $^\circ$C on 10-16) with
Eq.~(3) in Wexler's form; $\rho_i$ in the latent term; ``Allen-Cahn'';
``Modeling the''; the comma in Ahmed Mahjoub's name; \texttt{finalold}
class. The entries below list only what is still open.

\section*{Work list, in order}
\begin{enumerate}
\item \textbf{Correct Section~5.1 as drafted} (L426--477). Two sentences
  attribute the shortfall to surface diffusion, which Appendix~A rules out;
  ``by 41\% and 170\%'' misstates the numbers; one sentence is unfinished.
  Details under Section~5.1 below.
\item \textbf{Restate the gap against Bouvet et al.\ (2022)} (Intro \P8, F1)
  and fix the exponent sources (\P5, F4).
\item \textbf{Correct the $\alpha_c$ regime} in Intro \P6 (``diffusion
  limited'') and Section~3.2 (tex L455, ``vapor-limited''): $\alpha_c =
  10^{-3}$ is attachment-limited (F2).
\item \textbf{Rewrite Sections 3.2, 5.1 and Appendix~B for the tangent-start
  run} (plans and suggested text under each), now in Figure~4: the draft text (tex L455--461, L502) still describes
  the pre-necked run (a 28~$\mu$m starting neck, exponent 0.21). \Have \textbf{Tangent-start result (10-15)}: free exponent $a = 0.289$
(Kuczynski $1/m = 0.290$; 0.277 from later window starts), inside Demmenie's
0.26--0.33 and on the sharp-interface 0.29; no trapped pocket; neck to
61.2~$\mu$m in 150~h. The pre-necked 0.21 came from the initial condition.
Figure~4 is rebuilt from this run (\texttt{studies/molaro\_2019/demmenie/README.md}).
\item \textbf{Write Sections 2.4 and 2.5}; move the set-up text out of
  Sections 3 and 4 into them.
\item \textbf{Add the numbers to Section~4} and correct the high-porosity
  paragraph (L382--400).
\item \textbf{Finish 5.2; write 5.3--5.6, the Conclusion, Appendices A and
  B.} Plans below; the two-page scaling note
  (\texttt{temperature\_scaling/}) is the source for 5.2.
\item \textbf{Front matter:} key points, abstract, plain-language summary.
\item \textbf{Small corrections} (last section of this outline).
\end{enumerate}

\section*{Runs pending}
\begin{itemize}
\item \textbf{Saturated pair from tangent contact}: done (job 4474214).
  \Have \textbf{Tangent-start result (10-15)}: free exponent $a = 0.289$
(Kuczynski $1/m = 0.290$; 0.277 from later window starts), inside Demmenie's
0.26--0.33 and on the sharp-interface 0.29; no trapped pocket; neck to
61.2~$\mu$m in 150~h. The pre-necked 0.21 came from the initial condition.
Figure~4 is rebuilt from this run (\texttt{studies/molaro\_2019/demmenie/README.md}).
\item \textbf{Molaro $-20\,^{\circ}$C from tangent contact} (prepared, not
  submitted). A 10-minute pilot (batch file
  \texttt{molaro\_tangent\_T-20\_pilot}) checks first for a trapped void at
  the neck, the artifact that led to the pre-made neck in August.
  \Have Pilot run 10-15: the enclosed air pocket forms again (at
  $r = 3.3~\mu$m by 60~s), at half the interface parameter. Not submitted.
\item The passages marked ``provisional'' below can now be written from the
  tangent-start run.
\end{itemize}

\section*{Section titles}
The descriptive headings are adopted in the draft. ``2 Modelling
Approach'': AGU uses American spelling, and the draft writes ``modeling''
elsewhere, so ``Modeling''. Still to change: ``2.2 Gibbs-Thomson Condition''
(Interface Kinetics and the Sintering Time Scale), ``2.3 Homogenizing for
Effective Thermal Conductivity'' (Effective Thermal Conductivity by
Homogenization), and capitals in ``5.1 What vapor transport alone
reproduces''. All optional.

\section*{Build check (LaTeX warnings)}

The project uploaded on 2026-10-15 (second upload, \texttt{finalold} class)
compiles here in 25 pages with \textbf{no warnings} other than the one
below, and BibTeX reports none. What remains:
\begin{enumerate}
\item \textbf{``Label(s) may have changed. Rerun''} appears after any edit
  that moves a label and goes away on the next compile. Overleaf reruns
  automatically; if it persists there, recompile from scratch (the
  ``Recompile'' menu, ``Clear cached files'').
\item \textbf{Not cited in the text} (no warning, but AGU requires it):
  Table~2 (temperature-dependent quantities), Table~3 (complete
  parameters), Figure~2.
\item \textbf{Unused files} in the project: the Figure~8 timescale map,
  the all-porosity variant of Figure~7, and the AGU template file.
\end{enumerate}
If Overleaf shows warnings other than these, paste them: its compiler
settings can differ from the build here.

\section*{Front matter}
\begin{enumerate}
\Pg{Title and authors}{L1--4}{The paper's name and author list.}
  \begin{itemize}
  \item \Have ``Modeling the''; ``Ahmed Mahjoub''.
  \item \tG Remove ``--Version 2'' from the title.
  \end{itemize}
\Pg{Key points}{L7--8, placeholder}{Three statements of at most 140 characters each.}
  \begin{itemize}
  \item \tA For a given ice-grain packing, temperature from $-40$ to
    $-5\,^{\circ}$C only rescales the time over which sintering raises
    conductivity.
  \item \tA Effective conductivity follows specific surface area by one power
    law in well-connected packings; porosity sets its level.
  \item \tA At saturation, vapor transport alone gives the neck-growth
    exponent measured between ice spheres (0.29 against 0.26--0.33).
  \end{itemize}
\Pg{Abstract}{L9--31, marked ``needs to be updated''}{Background, method, the two-grain test, the aggregate result.}
  \begin{itemize}
  \item \tC ``Driven by temperature gradients'': every run here is
    isothermal.
  \item \tC ``replicating experimental results for various temperatures
    without retuning'' and ``incorporating a temperature dependence into our
    kinetic parameters measured in experiments'': neither is done;
    $\alpha_c$ is constant and the model gives 41\% of the measured growth.
  \item \tC ``collapse \ldots\ for both varying temperatures and varying
    porosities'': across temperature yes; across porosity only the exponent
    is shared, for $\varphi \le 0.375$.
  \item \tT The first five sentences are background; keep two.
  \item \tA Sentence by sentence: setting; gap (conductivity as grains sinter
    without compaction); method; grain pairs (41\% at $-20\,^{\circ}$C;
    exponent 0.29 at saturation, within the measured 0.26--0.33);
    aggregates (125 runs); result~1 (same
    conductivity against SSA at every temperature within 0.3\%); result~2
    (SSA$^{-0.80}$ for $\varphi \le 0.375$, level set by porosity); meaning.
    At most 250 words (it is 254 now).
  \end{itemize}
\Pg{Plain Language Summary}{L32--55, marked ``needs to be updated''}{The same story for a general reader.}
  \begin{itemize}
  \item \tC ``tune our model at $-20\,^{\circ}$C \ldots\ without retuning
    \ldots\ strong agreement'' and ``regardless of temperature or initial
    porosity'': as for the abstract.
  \item \tC ``changes in temperatures (sunrise or sunset)'' and ``gradients
    in temperature'': not part of this study.
  \item \tG ``know as'' (known); ``effecting'', ``effects'' twice, ``effect''
    (affecting, affects, affect); ``C.Then'' (space).
  \item \tT About 300 words; the limit is 200.
  \end{itemize}
\end{enumerate}

\section{Introduction}
\textbf{Re-synced 2026-10-15 against the current draft} (upload of 10-15,
source lines ``tex L''). Nine paragraphs. Each entry lists what the 10-15
draft already did (\Have), then only what is still open, with the
literature-check findings (F$n$, top of this document) merged in place.
\begin{enumerate}
\Pg{1. What dry snow metamorphism is}{tex L254}{Defines the process and says it changes connectivity and bulk properties.}
  \begin{itemize}
  \item No changes. Porosity stays in the list (metamorphism in general
    changes it; the model section says ours does not); the driving-force
    sentence is right as written; freeze/melt stays.
  \end{itemize}
\Pg{2. Why it matters on Earth}{tex L257}{Everyday examples, avalanches, and heat flow through snowpacks and firn.}
  \begin{itemize}
  \item \Have ``DSM is why''; the permafrost clause is gone; the
    conductivity sentence now cites Calonne et al.\ (2011, 2019) for heat
    flow and the temperature profile (supported: Calonne et al.\ 2019, p.2).
  \item \tC \Chk{F10} Depth-hoar sentence: Colbeck (1998) is a theory of bond
    growth (p.2: ``The basic shape of bonds in snow is dictated by the
    geometrical requirements of grain-boundary grooves''), not depth hoar or
    avalanches; drop it and keep Fukuzawa \& Akitaya and Schweizer et al.
    ``drives slab avalanche initiation'' $\to$ ``forms a weak layer on which
    slab avalanches can release''.
  \item \tC The Calonne citation is written \texttt{\textbackslash citet}
    mid-sentence; it must be \texttt{\textbackslash citep}.
  \item \tG ``glacial systems'': metamorphism acts in snow and firn, not in
    glacier ice; ``snowpacks and firn''. ``viability of ice in our alpine and
    polar regions'': ``persistence of ice in alpine and polar regions''.
    ``understanding for how'': ``understanding of how''.
  \item \tA (optional) Conductivity is also an input of detailed snowpack
    models, which set it from density alone (Vionnet et al.\ 2012, p.9,
    after Yen 1981). This sets up why a microstructural law matters.
  \end{itemize}
\Pg{3. Icy bodies and the Enceladus plume}{tex L260}{Water ice in the solar system, ocean worlds, the plume and what its grains are made of.}
  \begin{itemize}
  \item \Have The suggested wording is in (``Many bodies \ldots'',
    \textit{Cassini} italic, ``some of which is deposited back'').
  \item \tC \Chk{F10} ``contain percent level NaCl'': only a minority of
    grains do. Postberg et al.\ (2009, p.1): ``About 6\% of these spectra
    \ldots\ particularly high sodium contents''. Write ``some carry
    percent-level sodium salts, mainly NaCl''.
  \end{itemize}
\Pg{4. Sintering of plume deposits, and why conductivity}{tex L263, new on 10-15}{Where on Enceladus deposits sinter, how fast against burial, and why their conductivity matters.}
  \begin{itemize}
  \item \Have Drafted along the plan: Choukroun's experiment, the Tiger
    Stripes, burial, Byrne et al.'s conductivity range and why it matters.
  \item \tC \Chk{F6} Numbers and wording, all checked against the PDFs:
    \begin{itemize}
    \item ``may reach up to 180~K'': Spencer \& Nimmo (2013, p.8): ``surface
      temperatures within the troughs reach at least 180~K over regions tens
      of meters wide''.
    \item ``little to no sintering is expected'': Choukroun et al.\ (p.7)
      predict ``little to no strengthening''; they measure strength.
    \item ``0.01---1.0~mm\,yr$^{-1}$'' (Kempf; Southworth): the range is
      Choukroun et al.'s summary (p.7: ``on order of 0.01 to 1~mm/yr (Kempf et
      al., 2010; Southworth et al., 2019)''); Kempf et al.\ give ``0.5~mm per
      year in the immediate vicinity of the jet sources'' (p.10). Cite
      Choukroun with them. The rates assume compact ice (Southworth et al.,
      p.3: ``a lower bound on deposition'').
    \item ``so fresh deposits bury grains as they sinter'' needs Choukroun et
      al.\ (p.7: ``competition between strengthening of existing plume
      deposits and covering by new fresh and loose particles'').
    \item ``Sintering that strengthens a deposit will also cause changes to
      \ldots\ thermal conductivity'' needs a citation: Molaro et al.\ (p.23):
      neck growth causes ``an increase in the thermal conductivity of the
      aggregate''.
    \item Three \texttt{\textbackslash citet} that should be
      \texttt{\textbackslash citep}; ``preformed'', ``grows linear'', ``with
      about 15 years'', ``expected be greatest''; en dash in the range.
    \end{itemize}
  \item \tC \Chk{F11} \texttt{byrneRadarAttenuationSubsurface2024} is the
    preprint; use the JGR Planets article, doi:10.1029/2024JE008615. Their
    conductivities are for an airless surface; ours are with air at 1~atm,
    so do not place our values on their scale.
  \Snippet{Under typical Enceladus surface conditions, however, little to no
    strengthening is expected, and the deposits remain poorly consolidated
    granular packings \citep{choukroun_strength_2020}. Near the Tiger
    Stripes, surface temperatures reach at least 180~K
    \citep{spencerEnceladusActiveIce2013}, and there the deposits could
    consolidate within about 15 years \citep{choukroun_strength_2020}.
    These are also the regions of heaviest plume deposition, on the order of
    0.01--1~mm\,yr$^{-1}$ \citep{kempfHowEnceladusDust2010,
    southworthSurfaceDepositionEnceladus2019, choukroun_strength_2020}, so
    fresh grains bury a deposit while it sinters
    \citep{choukroun_strength_2020}. As the necks grow, the conductivity of
    the deposit rises \citep{molaroMicrostructuralEvolutionWater2019}.}
  \end{itemize}
\Pg{5. Sintering as part of metamorphism}{tex L266; Figure~1 at L271}{Neck growth between grains, which transport route dominates, earlier work.}
  \begin{itemize}
  \item \Have Choukroun is out of the vapor-dominance sentence; the
    snow-metamorphism examples are gone; the exponent sentence is in, with
    Hobbs \& Mason, Kingery, Thomas and Demmenie in the \texttt{.bib}; Molaro
    et al.'s grain-pair measurements are named; Figure~1's caption is
    vapor-only.
  \item \tC \Chk{F4} \textbf{The exponent sentence cites the wrong sources.}
    Hobbs \& Mason's model, with vapor diffusing through air to the neck,
    gives $1/5$, not $1/3$ (Blackford 2007, p.16: ``Hobbs and Mason's modified
    model gave values of $n = 5$ and $m = 3$ for vapour transport''; their
    PDF is not in the library). The $1/3$, $1/5$, $1/7$ set is Kingery (1960,
    Table~I) and Demmenie et al.\ (2025, Table~1). Thomas et al.\ give ``$n=3$
    for vapour transport, $4\le n\le5$ for surface diffusion'' (p.3); cite
    them for what they measured: ``a neck growth proportional to $t^{1/3}$,
    i.e., the sublimation and recondensation process predominate at
    $-20\,^{\circ}$C'' (p.4).
  \item \tC \Chk{F4} ``$n$ identifies the dominating mechanism'': Maeno \&
    Ebinuma (1983, p.4) call that ``misleading because the sintering rate is
    usually the sum of contributions from several mechanisms''. And name the
    debate: Kingery measured ``$(x/r)\sim t^{1/6.9}$, in close agreement with
    \ldots\ surface diffusion'' (p.5); Demmenie et al.\ (p.2) attribute it to
    ``an underestimation of the relative humidity''.
  \item \tC \Chk{F4} ``Other modes \ldots\ operate over planetary timescales
    (Molaro)'' reverses Molaro et al.\ (p.30): at 200~K neglecting surface
    diffusion changes their $\tau_{50}$ by ``$\sim$3,700\%'', and they
    ``recommend always including the surface diffusion term \ldots\ in
    planetary applications''. At 250~K it matters only for the smallest
    grains (13\% for 1~$\mu$m; ``no significant effects'' for 100~$\mu$m).
    Drop ``or operate over planetary timescales''; this also supports not
    extrapolating to Enceladus temperatures.
  \item \tC The vapor-dominance sentence has no citation: Thomas et al.\
    (1994, p.4) and Demmenie et al.\ (2025) are the direct ice
    measurements.
  \item \tA \textbf{What Molaro et al.\ measured} is still thin: neck growth
    between grain pairs in a sealed cryostage at $-20$ and $-5\,^{\circ}$C.
    They report no exponent; $0.23\pm0.03$ is our fit to their points.
  \item \tA \textbf{Describe Demmenie et al.'s experiment in one sentence.} In
    \P5 they are only one key in a list of four citations, but Section~3.2
    tests the model against their numbers, so the reader needs to know what
    they measured: two ice spheres brought into contact in a box held at
    ice saturation, at $-3\,^{\circ}$C, with exponents 0.26--0.33 (p.3), and
    ``$\sim$1/7'' when ``we deliberately set the humidity too low and let the
    ice sublimate'' (p.3).
  \item \tA \textbf{Scope: rounded grains, not fresh snow} (discussed 10-16).
    Molaro et al.\ (spheres sprayed into liquid nitrogen, p.15), Demmenie et
    al.\ (ice spheres) and our aggregates (discs in point contact) all start
    from rounded grains. None of them is snow, fresh or aged, so they cannot
    be said to use ``snow that has already metamorphosed''. What can be
    said: fresh snow first rounds, and that stage is fast; Kaempfer \&
    Schneebeli (2007) saw ``the transition from the dendritic form of the
    fresh snow at time zero to rounded grains a month later'' (p.5), the
    first phase of ``a two-stage process'' (p.1). Sphere experiments and our
    model describe the second stage, sintering of rounded grains. Our
    results do not cover the first stage, so they cannot claim it. Put this
    in Section~2.1 (Scope) and Section~5.6.
  \item \tG ``an adapted model sintering model originally used in metallurgy
    first proposed by'': Molaro et al.\ made ``No modifications'' to it (p.5),
    and Swinkels \& Ashby's is a ``second report''; ``applied the sintering
    model of \citet{swinkelsSecondReportSintering1981}, developed for
    metals''. ``Sintering has been studied extensively \ldots'' has no
    citation: Blackford (2007) is a review of ice sintering. The last
    sentence (Choukroun) repeats \P4: cut.
  \item \tG ``In a laboratory setting moves tends be redistributed''; ``Neck
    growth tend s''; $w$ is ``the neck width'' here but ``a neck of radius
    $w(t)$'' in the Figure~1 caption: width throughout. Caption: ``fo''.
  \Snippet{In laboratory experiments, mass reaches the neck mainly by vapor
    transport: ice sublimates away from the neck and desublimates on it
    \citep{thomas_crushing_1994, demmenieIceSinteringSublimation2025}, as
    shown in \cref{fig:Figure1_mechanisms_aggregate}. At these temperatures,
    other routes matter only for the smallest grains, but surface diffusion
    gains weight as temperature falls
    \citep{molaroMicrostructuralEvolutionWater2019}. Neck growth often
    follows $w \propto t^{n}$. When attachment at the ice surface sets the
    rate, vapor transport gives $n=1/3$, while volume and surface diffusion
    give $1/5$ and $1/7$ \citep{kingery_regelation_1960,
    demmenieIceSinteringSublimation2025}; when vapor must diffuse through
    air to reach the neck, vapor transport also gives $n=1/5$
    \citep{blackfordSinteringMicrostructureIce2007}. Several mechanisms act
    at once, so the exponent alone does not identify the dominant one
    \citep{maenoPressureSinteringIce1983}. \citet{kingery_regelation_1960}
    measured $n \approx 1/7$ and attributed it to surface diffusion, whereas
    \citet{demmenieIceSinteringSublimation2025} measured $n$ between 0.26
    and 0.33 at saturation and close to $1/7$ when the humidity was too low
    and the ice sublimated, and attributed Kingery's value to an
    underestimated humidity.}
  \end{itemize}
\Place{Figure 1}{\Have placed. Caption: see ``Figure captions'' after Section~6.}
\Pg{6. Attachment kinetics}{tex L276}{Supersaturation, the Hertz--Knudsen law, and the condensation coefficient.}
  \begin{itemize}
  \item \Have Libbrecht \& Rickerby (2013) now cited for the measurements, on
    basal and prism facets; the layer-nucleation wording; the form of
    $\alpha_c$; the range cited; a closing sentence on the constant
    $\alpha_c$; ``The growth rate of ice from vapor depends on''.
  \item \tC \Chk{F2} \textbf{``$\alpha_c=10^{-3}$ to force a diffusion limited
    growth regime'' is backwards.} A smaller $\alpha_c$ slows attachment and
    moves growth toward attachment control. Bouvet et al.\ (2024, p.5) place
    the transition at ``$10^{-3}$ for small grains ($\sim$0.1~mm)'', about the
    diameter of our grains ($\bar R = 50~\mu$m). The draft's Section~5.1
    (tex L502) already says ``attachment limited regime''; Section~3.2 (tex
    L455) says ``vapor-limited'' and needs the same fix.
  \item \tC \Chk{F8} Libbrecht \& Rickerby do not state the range
    $10^{-3}$--$10^{-1}$; Braun et al.\ (2024, p.12) call it ``the commonly
    found range'', ``also used by Kaempfer and Plapp (2009)''. Cite those two;
    ``bounded'' $\to$ ``commonly used''. $\alpha_c=\exp(-\sigma_0/\sigma)$ is not
    an Arrhenius function (its argument is $1/\sigma$, not $1/T$); only
    $\sigma_0$ depends on temperature; ``surface saturation'' $\to$
    ``supersaturation''.
  \item \tG ``led an experiment in near vacuum'': ``measured at reduced
    pressure''. ``and thus slowing attachment'': ``which slows attachment''.
  \Snippet{Values of $\alphaC$ between $10^{-3}$ and $10^{-1}$ are commonly
    used \citep{kaempferPhasefieldModelingDry2009,
    braunRigorousApproachSpecific2024}. In the layer-nucleation model of
    \citet{libbrechtMeasurementsSurfaceAttachment2013},
    $\alphaC = A\exp(-\sigma_0/\sigma_\mathrm{surf})$, where the critical
    supersaturation $\sigma_0$ depends on temperature and differs between the
    basal and prism facets. In this study we hold $\alphaC$ constant. We use
    $\alphaC=10^{-1}$, the upper end of the range, to compare with the
    experiments of \citet{molaroMicrostructuralEvolutionWater2019}, and
    $\alphaC=10^{-3}$ for the saturated grain pair and the aggregates. For
    grains of about 0.1~mm, $\alphaC=10^{-3}$ lies near the transition
    between attachment-limited and diffusion-limited growth
    \citep{bouvetMultiscaleModelingHeat2024}.}
  \end{itemize}
\Pg{7. Metamorphism beyond a pair; earlier models}{tex L279}{From a grain pair to an aggregate, and the models that exist.}
  \begin{itemize}
  \item \Have Braun et al.\ described and the $\alpha_c$ fits stated as
    suggested; Libbrecht (2017) out of the DSM list; ``all stages of DSM'';
    ``to better predict''.
  \item \tC \Chk{F10} ``Many experiments attempt to parameterize'' cites
    Colbeck (1983), a theory paper, and Libbrecht (2017), a review: cite
    Libbrecht \& Rickerby (2013).
  \item \tA (optional) End on Braun et al.'s reading, which is the evidence
    for our own use of $\alpha_c$: they call it an ``effective condensation
    coefficient'' (p.13).
  \item \tG ``x-ray compute microtomography'': ``X-ray computed
    microtomography''. ``were the first to attempt to use'': ``were the first
    to use''. ``steady state'' (adjective): ``steady-state''.
  \end{itemize}
\Pg{8. What is not known}{tex L282}{The gap the paper addresses.}
  \begin{itemize}
  \item \Have ``affects''; ``bulk properties of porous ice''; the paragraph
    now ends on the paper's own gap; Calonne et al.\ (2011) and a gradient
    sentence added.
  \item \tC \Chk{F1} \textbf{The gap is partly closed already.} ``How
    isothermal metamorphism influences the macroscopic properties in the
    absence of compaction is still unquantified'' is contradicted by
    Bouvet et al.\ (2022), cited in this paragraph: they computed
    $k_{\mathrm{eff}}$ on snow evolved isothermally by their phase-field model
    (p.12: ``conductivity increases from 0.078 to 0.081~W\,m$^{-1}$\,K$^{-1}$'';
    ``changes are rather small''), and could not separate microstructure
    from density (p.15: ``the impact of the microstructural changes is not
    visible''). What remains is ours: grains sintering from point contacts
    at fixed porosity, conductivity against SSA, and the effect of
    temperature and porosity.
  \item \tC \Chk{F1} The gradient sentence follows Calonne et al.\ (2011), a
    survey of 30 samples with no time series. Its evidence is Calonne et
    al.\ (2014): SSA fell ``from 27.7 to 13.4~m$^2$\,kg$^{-1}$'' (p.9);
    conductivity showed ``a decrease between 0 and 73~h, followed by an
    increase'' (p.16), ``an overall increase \ldots\ at a constant density''
    in their words. Schneebeli \& Sokratov (2004, p.10) saw a rise ``by about
    60\%''. Our reconstruction from their Table~1 (below) has the
    orientation average ending near its start while $k_z$ rose; say that the
    loss of SSA did not by itself raise conductivity, not that it never did.
  \item \tC \Chk{F10} ``centered around the snow community here on earth''
    cites Flin et al.\ (2003, 2004), who report microstructure, not bulk
    properties: cite Calonne et al.\ (2011, 2014). ``here on earth'': ``on
    Earth''.
  \item \tA \textbf{Why SSA}, one sentence: SSA is measured in the field and
    laboratory (Matzl \& Schneebeli 2006; Gallet et al.\ 2009; Proksch et
    al.\ 2015), so a law in SSA can be applied. Still missing.
  \item \tA \textbf{What the conductivity here includes}, one clause: under a
    gradient, vapor transport adds to the apparent conductivity (Fourteau et
    al.\ 2021); ours is conductive only, which is exact for isothermal
    sintering. Still missing (here or in Section~2.3).
  \item \Note \textbf{Background on snow conductivity} (for the two
    sentences above; from the earlier outline). Calonne et al.\ (2011):
    full conductivity tensor from tomography of 30 samples, a regression on
    density ($r = 0.985$), scatter from anisotropy ($k_z/k_{xy}$ 0.7--1.5).
    Calonne et al.\ (2014): rounded grains, 294~kg\,m$^{-3}$, 43~K\,m$^{-1}$
    at $-4\,^\circ$C for 3 weeks; density constant; SSA 27.7 $\to$ 13.4;
    density-corrected $k_z$ 0.18 $\to$ 0.17 $\to$ 0.23, $k_{xy}$ 0.20 $\to$
    0.15 $\to$ 0.18 (our reconstruction, $\pm$0.01); $k_z/k_{xy}$ 0.90
    $\to$ 1.34; they predict $k$ from porosity and correlation lengths, never
    from SSA. Calonne et al.\ (2019): $k_z/k_{xy}$ tracks the
    correlation-length ratio; at low density ``the few bonds between ice
    grains \ldots\ act as heat paths'' (p.5).
  \Snippet{Under a temperature gradient, \citet{calonneStudyTemperatureGradient2014}
    followed a snow layer at constant density for three weeks: its SSA
    halved, while its conductivity first decreased, as bonds between grains
    sublimated, and then increased, mostly in the vertical direction, a
    two-stage change that \citet{schneebeliTomographyTemperatureGradient2004}
    had also seen. Losing surface area did not by itself raise the
    conductivity. Under isothermal conditions,
    \citet{bouvetSnowEquiTemperatureMetamorphism2022} computed the
    conductivity of snow evolved by a phase-field model and found only small
    changes, which followed small variations in density. Neither study
    follows a packing of ice grains as it sinters from point contacts at
    fixed porosity, relates the conductivity to the SSA, or asks how
    temperature and porosity change that relation.}
  \end{itemize}
\Pg{9. What follows}{tex L285}{The order of the paper and its findings.}
  \begin{itemize}
  \item \Have The two findings are rewritten: the conductivity--SSA relation
    is independent of temperature, and for $\varphi \le 0.375$ it is a power
    law with a shared exponent.
  \item \tC Porosity: say also that the level falls with porosity (about
    threefold from $\varphi = 0.275$ to 0.475); ``same exponent'' alone
    hides it.
  \item \tC ``an updated phase-field model introduced by Moure \& Fu'': Moure
    \& Fu (2024) is a wet-snow model; ``the ice--air form of the phase-field
    model of Moure \& Fu (2024)''.
  \item \tC ``use their results to compare with their model'': ``compare our
    model with their measurements''.
  \item \tC ``a diffuse interface adapted homogenization technique initially
    outlined by Auriault (1991)'': ``a homogenization method
    \citep{auriaultHeterogeneousMediumEquivalent1991} adapted to the diffuse
    interface''.
  \item \tA \textbf{Announce the two-grain results} in one sentence: vapor
    transport alone gives 41\% of the neck growth Molaro et al.\ measured at
    $-20\,^{\circ}$C, and the exponent Demmenie et al.\ measured at
    saturation (0.29 against 0.26--0.33).
  \item \tA \textbf{The temperature result is not announced}: temperature
    from $-40$ to $-5\,^{\circ}$C only rescales time (Section~5.2). One
    clause.
  \item \tA \textbf{Scope, in one clause}: isothermal, air at 1~atm, two
    dimensions, temperatures a laboratory can reach.
  \item \tA \textbf{What is new}, one clause: the thin-interface kinetics and
    the homogenization in the diffuse interface (tensor law; prior art
    Nicoli et al.\ 2011).
  \item \tG ``In what follows'' opens this paragraph and the last sentence of
    \P8: once. ``we will present'': ``we present''. ``over 2D packings
    \ldots\ over 30-day time periods'': ``in two-dimensional packings of ice
    grains for 30 days''. ``tracks the temperature and water vapor density
    fields'': and the phase field.
  \end{itemize}
\end{enumerate}

\section{Modeling Approach}
\textit{Draft heading: 2 Modelling Approach.}
\begin{enumerate}
\Pg{Section opening}{tex L293}{Says the model follows Kaempfer \& Plapp and Moure \& Fu.}
  \begin{itemize}
  \item \Have The novelty sentence (``The contribution here is the
    application of the phase-field model, not the numerical methods''); the
    homogenization lineage (Auriault 1991; Calonne et al.\ 2011, 2014, 2019).
  \item \tA \Chk{F10} Calonne et al.\ (2014, p.6) cite the Auriault et al.\
    (2009) book for this boundary-value problem; add it next to Auriault
    (1991). The book is in the library; it is not in the \texttt{.bib}.
  \item \tA One sentence of lineage from the Discussion's second paragraph
    can move here.
  \end{itemize}
\subsection{Phase-Field Model of Dry Snow Metamorphism}
\Pg{What a phase field is}{tex L297}{General description of the method.}
  \begin{itemize}
  \item \tG ``representing the volume fraction of phase $l$ at a given point
    \ldots\ and evolves through time'': the sentence changes subject; ``which
    evolves in time''.
  \item \tT Textbook material; one sentence.
  \end{itemize}
\Pg{Unknowns}{tex L300}{Lists the four fields.}
  \begin{itemize}
  \item \Have $\phi_a$ is ``air''.
  \item \tC \textbf{Temperature in $^\circ$C (user, 10-16).} $T$ is in
    $^\circ$C, which is how every result is reported; the solver works the
    same way (\texttt{RhoVS\_I}: \texttt{temK = tem + 273.15}). Introduce the
    absolute temperature $T + T_0$ ($T_0 = 273.15$~K) once, here, and use it
    only where a formula needs kelvin: Eq.~(3) (Wexler), Eq.~(5) ($D_v$),
    Eq.~(7) ($\beta_{\mathrm{sub}}$, $k_B (T + T_0)$) and $d_0$. The
    conservation equations (1b), (1c) and (11) involve only $\partial T$ and
    $\nabla T$, so they are unchanged. Table~1 [written] 10-16.
  \Snippet{The temperature $T$ is expressed in degrees Celsius; where a
    relation requires the absolute temperature we write $T + T_0$, with $T_0 =
    273.15$~K.}
  \item \tG ``derived from the free energy potential and allows'': ``that
    allows''.
  \item \tT The last two sentences (``defined throughout the domain'',
    ``benefit from the phase-field formulation'') can go.
  \end{itemize}
\Pg{Scope}{to write}{Says which part of dry snow metamorphism the model contains.}
  \begin{itemize}
  \item \tA The model carries sintering by vapor transport: sublimation,
    diffusion through the pore, desublimation.
  \item \tA Not in it: surface diffusion, grain-boundary processes,
    densification, overburden. Define desublimation.
  \item \tA Vapor-transport sintering does not densify: grain centres do not
    approach, so porosity is constant in every run.
  \item \tA Hence the laboratory analogues are unloaded ice-sphere samples,
    not natural snow.
  \end{itemize}
\Pg{Model equations and the phase equation}{tex L304--311}{Eq.~(1) and the meaning of its terms.}
  \begin{itemize}
  \item \Have Eq.~(1b) carries $\rho_i$; ``$\varepsilon$ controls the
    width''; ``Allen-Cahn''.
  \item \tC \Chk{F9} $M_{\mathrm{sub}}$ and $\alpha_{\mathrm{sub}}$ now point
    to the appendix ``Model Scaling Parameters'' (tex L527), which is an
    empty heading; label typo \texttt{Sclaing}. Content:
    $M_{\mathrm{sub}} = \varepsilon/(3\tau_{\mathrm{sub}})$,
    $\alpha_{\mathrm{sub}}$, $\beta_{\mathrm{sub}}$, $\Delta\beta$,
    $\tau_{\mathrm{sub}}$ (now in the notes to Table~1).
  \item \tC Eq.~(1c) as printed has no $\xi_v$, but Tables~3 and~4 list
    $\xi_v = 10^{-3}$. Show it or explain it in one sentence.
  \item \tG ``The product \ldots\ localize this phenomena'': ``localizes this
    term''.
  \end{itemize}
\Pg{Saturation vapor density}{tex L313--319}{Eqs.~(2) and (3).}
  \begin{itemize}
  \item \Have Eq.~(3) is now Wexler's (1977) Eq.~54, as used by
    Kaempfer \& Plapp (2009); ``over ice''; ``are''.
  \item \tC \Chk{F9} Eq.~(2) still writes $P_a$; the text says
    $P_{\mathrm{atm}}$. ``where $T_0 = 273.15$~K'' follows Eq.~(3), which
    does not contain $T_0$ yet. With $T$ in $^\circ$C, Eq.~(3) reads
    $P_{vs}^{I} = \exp\big(\sum_{j=0}^{4} k_j (T + T_0)^{j-1} + k_5 \ln
    (T + T_0)\big)$, and the sentence then defines $T_0$ (snippet below).
  \item \tC \Chk{F9} \textbf{Table [written] 10-16:} Eq.~(3) is Wexler
    (1977), Eq.~54 (p.11), coefficients $K_j$ in his Table~2 (p.9); the
    solver's $k_j$ are these rounded to four figures ($+0.04$ to $+0.05$\% in
    $P_{vs}$; Wexler's Table~3 gives 103.276~Pa at $-20\,^\circ$C). Table~1's
    note b now cites Wexler Eq.~54 and lists the six values; ``ASHRAE'' is
    gone (ASHRAE uses Hyland \& Wexler 1983, a different fit).
    \textbf{Open (text):} the sentence after Eq.~(3).
  \Snippet{where $T_0 = 273.15$~K and $P_{vs}^{I}$ in pascals
    \citep{wexlerVaporPressureFormulation1977}; the coefficients $k_j$ are
    listed in \cref{tab:material_properties}.}
  \end{itemize}
\Pg{Energy equation}{tex L321}{Eq.~(1b) and the mixture rules.}
  \begin{itemize}
  \item \tC ``latent heat released and absorbed during sublimation and
    desublimation, respectively'': reversed. Sublimation absorbs heat;
    desublimation releases it.
  \item \tG ``phase dependent'': hyphenate.
  \end{itemize}
\Pg{Vapor equation}{tex L331--335}{Eq.~(1c) and the diffusivity.}
  \begin{itemize}
  \item \tC With $T$ in $^\circ$C, Eq.~(5) becomes $D_v = D_{v0}\,\big((T + T_0)/T_0\big)^{1.81}$,
    with $D_{v0}$ the diffusivity at 0\,$^\circ$C. Tables~2--4 used $T_0$ for
    the run temperature in $^{\circ}$C; since 10-16 they write $T$, so $T_0$
    is only 273.15~K.
  \item \tG ``The rest of the thermophysical properties are shown in'':
    ``Material properties are listed in''.
  \item \tA Massman (1998) for $D_v(T)$, as Moure \& Fu do; not in the
    library, so not checked here.
  \end{itemize}
\Place{Table 1}{\Have Material properties, cited at tex L335.}
\end{enumerate}

\subsection{Interface Kinetics and the Sintering Time Scale}
\textit{Draft heading: 2.2 Gibbs-Thomson Condition.}
\begin{enumerate}
\Pg{The interface condition}{L216--226}{Eqs.~(6) and (7).}
  \begin{itemize}
  \item \tG ``is commonly interchanged with the condensation coefficient'':
    ``is commonly expressed through''. ``molecular mass a water molecule''
    (mass of a water molecule).
  \item \tA $d_0 = \gamma m/(\rho_{\mathrm{ice}} k_B T)$ is defined only in
    Table~1.
  \end{itemize}
\Pg{The thin-interface limit}{to write}{What ``under the conditions outlined by Kaempfer \& Plapp'' means.}
  \begin{itemize}
  \item \tC \textbf{Checked 10-16 (user asked whether the correction needs
    isothermal conditions).} The literature condition is not isothermality.
    \begin{itemize}
    \item The thin-interface correction $\beta' = a_1(\tau/W\lambda - a_2
      W/D)$ (Kaempfer \& Plapp 2009, Eq.~39, p.8, from Karma \& Rappel) is
      derived for \emph{one} diffusing field with ``the diffusion
      coefficient $D$ \ldots\ constant and equal in both phases'' (p.8),
      ``the standard symmetric model (with equal diffusivity in both
      phases)'' (Karma 2001, p.1). In that model the field is the
      temperature, so it is not an isothermal result.
    \item Our model has two coupled fields, and vapor diffuses only in the
      air (one-sided). For asymmetric diffusion a finite $W$ ``magnifies
      several nonequilibrium effects at the interface'' (a jump across it,
      interface stretching, surface diffusion; p.8), removed only with an
      antitrapping current. The thermosolutal version (Ramirez et al., their
      Eq.~41) has a coefficient $a_3$ that ``depends on the local interface
      temperature''.
    \item Kaempfer \& Plapp did not derive the thin-interface limit for
      their snow model: ``A simpler way \ldots\ is to observe that from
      simple dimensional arguments, the thin-interface corrections \ldots\
      must scale as $W/\alpha_T$ and $W/D_v$'' (p.8). They require these to
      be small, Eqs.~(42)--(43), so the \emph{sharp}-interface limit holds,
      plus $W \ll 1/\mathcal{K}$ and $\tau v_n/W \ll 1$ (p.9); ``a proper
      thin-interface limit will have to be developed'' (p.9). Moure \& Fu
      (2024, p.6): the asymptotic analyses ``do not readily apply'' with
      asymmetric diffusivities; antitrapping solutions exist ``only with
      symmetric diffusivities''; they impose $\varepsilon < D\beta'$,
      Eqs.~(27)--(29).
    \end{itemize}
    So $\Delta\beta = a_1 a_2\,\varepsilon\,(D_T^{-1} + D_v^{-1})\,
    \rho_{\mathrm{ice}}/\rho_{vs}^{I}$ is an \emph{estimate} built from that
    dimensional scaling with the symmetric-model $a_2$, not a derived limit
    of our equations. Where isothermality helps: with no imposed gradient
    $T$ varies only through latent heat, so the thermal part of the
    correction and the temperature dependence of $a_3$ hardly enter; the
    one-sided vapor part remains. Our own check: the lunar velocity study
    matched series theory to 1--3\% with the term on.
    \textbf{For the draft:} where the correction is 2\% (aggregates) or
    0.2\% (saturated pair), we are in Kaempfer \& Plapp's sharp-interface
    regime and the form of $\Delta\beta$ hardly matters; at 23\% (Molaro
    pairs) it does, and the text should say so.
  \Snippet{Our model reproduces Eq.~(6) in the sharp-interface limit, which
    requires the interface parameter to be small compared with the radius
    of curvature and with the lengths $D_v\beta_{\mathrm{sub}}\rho_{vs}^{I}/\rho_{\mathrm{ice}}$
    and $D_T\beta_{\mathrm{sub}}\rho_{vs}^{I}/\rho_{\mathrm{ice}}$
    \citep{kaempferPhasefieldModelingDry2009, mourePhaseFieldModelWet2024}.
    For a finite interface we add to $\beta_{\mathrm{sub}}$ the correction
    $\Delta\beta$ (\cref{tab:material_properties}), whose scaling follows
    \citet{kaempferPhasefieldModelingDry2009}; it has not been derived for
    a model with one-sided vapor diffusion, so we treat it as an estimate.
    It is 2\% of $\beta_{\mathrm{sub}}$ in the aggregate simulations and
    23\% in the simulations of the experiments of
    \citet{molaroMicrostructuralEvolutionWater2019}.}
  \item \tA Its size: 2\% of $\beta_{\mathrm{sub}}$ in the aggregates, 0.2\%
    in the saturated pair, 23\% in the Molaro pairs.
  \end{itemize}
  \Cite{Kaempfer \& Plapp 2009 \texttt{kaempferPhasefieldModelingDry2009}; Moure \& Fu 2024 \texttt{mourePhaseFieldModelWet2024}.}
\Pg{The time scale}{to write here; now at L469--477 in the Discussion}{Defines $\tau_{\mathrm{sub}}$ and the sintering age $\theta$.}
  \begin{itemize}
  \item \tA Move Eq.~($\tau_{\mathrm{sub}}$) from Section~5.1 to here, and
    define $\theta = t/\tau_{\mathrm{sub}}$ here, not in Section~4.2.
  \item \tA Two time scales of the same form $\ell^{2}\beta/d_0$:
    $\tau_{\mathrm{sub}}$ ($\ell = \varepsilon$), the relaxation time of
    the diffuse interface; and $\tau_{\bar R} = \bar R^{2}\beta/d_0$, the
    time for a grain to change shape. $\tau_{\mathrm{sub}} =
    \tau_{\bar R}/2500$ in every aggregate run.
  \item \tA The 30-day runs end at $\theta = 43$ ($-40\,^{\circ}$C) to 1256
    ($-5\,^{\circ}$C). Cite Table~2 here.
  \end{itemize}
\Pg{The values used}{to write}{The condensation coefficients and why they are constant.}
  \begin{itemize}
  \item \tA $\alpha_c = 1.020\times10^{-3}$ (aggregates), 0.1 (Molaro pairs),
    $1.002\times10^{-3}$ (saturated pair), constant in temperature.
  \item \tA Why constant: a temperature-dependent $\alpha_c$ would make a
    temperature sweep uninterpretable.
  \end{itemize}
  \Cite{range of $\alpha_c$: Braun et al.\ 2024 \texttt{braunRigorousApproachSpecific2024}; Bouvet et al.\ 2022 \texttt{bouvetSnowEquiTemperatureMetamorphism2022}; Libbrecht 2017 \texttt{libbrechtPhysicalDynamicsIce2017}.}
\Place{Table 2}{\Have Temperature-dependent quantities, placed. \tA Not cited
  in the text.}
\end{enumerate}

\subsection{Effective Thermal Conductivity by Homogenization}
\textit{Draft heading: 2.3 Homogenizing for Effective Thermal Conductivity.}
\begin{enumerate}
\Pg{Two scales and the microscale problem}{L228--238}{Scale separation, periodicity, Eq.~(8).}
  \begin{itemize}
  \item \tC ``without less of generality'': periodicity is an assumption,
    not a convenience (and the phrase is ``loss'').
  \item \tG ``A smaller length scale $l$ describing \ldots, and $L$
    describing \ldots'' is not a sentence. ``representative element
    volume'' (elementary). ``diffusive heat transfer'' (conductive).
    ``upscaling effective thermal conductivity'' (computing the).
  \end{itemize}
\Pg{Macroscale law and cell problem}{L239--247}{Eqs.~(9)--(11).}
  \begin{itemize}
  \item \tC \Chk{F9} Eq.~(11d) (tex L379) has $\mathbf{t}$ where it needs
    $\nabla\mathbf{t}$, as in Eqs.~(11a, b) and in Calonne et al.\ (2014,
    p.6, their Eq.~15).
  \item \tG ``are commonly referred to as co the solutions'': the sentence
    is broken. ``the corresponding macroscopic description is described
    as''.
  \item \tT ``Equations \ldots\ are then coupled with the boundary conditions
    \ldots'' says the equations go with their own boundary conditions; cut.
  \end{itemize}
\Pg{Single-domain form}{L248--260}{Weak form and the indicator function, Eqs.~(12)--(15).}
  \begin{itemize}
  \item \tC The test function is scalar $v$ in the text and bold
    $\mathbf{v}$ in Eq.~(12). $K_\star$ in Eqs.~(13)--(14) becomes $K_{*}$
    in Eq.~(15).
  \item \tG ``as diffuse interface model'' (as a); ``in 2.1'' (Section~2.1);
    ``We can the define'' (then); ``$\zeta = 1.0$'' (1).
  \item \tT Eqs.~(12)--(14) can be one sentence: summing the weak forms over
    the two subdomains, the interface terms cancel by flux continuity, and
    Eq.~(15) follows.
  \end{itemize}
\Pg{Diffuse form and the tensor law}{L261--275}{Eqs.~(16), (17) and why the conductivity is a tensor in the interface.}
  \begin{itemize}
  \item \tC ``interpolates between 0 and 1'': between $K_{\mathrm{air}}$ and
    $K_{\mathrm{ice}}$.
  \item \tG ``along the our phase-field''; ``$K_\parallel(\phi)$ and are
    $K_\perp(\phi)$ the''; ``We then may write diffuse BVP'' (the diffuse);
    ``tonsorial'' (tensorial); ``interface width $\varepsilon$''. $\phi$
    and $\phi_i$ are used for the same field; $t_\star$ is not bold.
  \item \tA Credit and what is new: Nicoli, Plapp \& Henry (2011) define the
    law, show the cancellation analytically for a planar interface, and test
    it on a disk (contrast 2) and two half-disks (contrast 10); here it is applied to the cell problem at a contrast of 115
    and verified on laminate and disk-array cases (supplement).
  \item \tC \Chk{F10} ``both of order $\varepsilon$ for a scalar conductivity
    law'' (tex L421): a scalar law can cancel one of the two interface
    effects but not both; Nicoli et al.\ (p.1): ``both effects can be
    eliminated simultaneously only if the mobility is allowed to become a
    tensor''. Write ``no scalar law removes both''.
  \item \tA Why it matters here: with a scalar law the bias reaches tens of
    percent and grows with $\varepsilon\,$SSA, so it would drift as the
    aggregate sinters.
  \item \tA The cell problem is solved on frozen snapshots of $\phi$.
  \end{itemize}
  \Cite{Nicoli, Plapp \& Henry 2011 \texttt{nicoli\_tensorial\_2011}; Auriault 1991 \texttt{auriaultHeterogeneousMediumEquivalent1991}; Calonne et al.\ 2011 \texttt{calonneNumericalExperimentalInvestigations2011}.}
\Pg{What is reported}{to write}{Defines the quantities the results plot.}
  \begin{itemize}
  \item \tA $\mathbf{K}_{\mathrm{eff}}$ in the diffuse form, the one
    evaluated (Eq.~10 is the sharp form only).
  \item \tA The scalar $k_{\mathrm{eff}} = (k_{xx} + k_{yy})/2$, now defined
    only in the Figure~6 caption.
  \item \tA SSA: how it is measured from $\phi$, and its units. It is used
    throughout Section~4 and defined nowhere.
  \end{itemize}
  \Cite{for computing SSA from a field, possibly Flin et al.\ 2014 \texttt{flinComputationsSpecificSurface2014} (check that the method matches).}
\Place{Figure 2}{\Have placed. \tA Not cited in the text. Caption: see ``Figure captions'' after Section~6.}
\end{enumerate}

\subsection{Numerical Implementation}
\begin{enumerate}
\Pg{Discretization and time stepping}{to write; heading only at L276}{Enough for someone to reproduce the runs.}
  \begin{itemize}
  \item \tA Quadratic $C^{1}$ splines; element size $\varepsilon/\sqrt2$.
  \item \tA Adaptive step capped at $2\,\tau_{\mathrm{sub}}$ in the
    aggregates. The grain-pair runs use larger caps (5~$\tau_{\mathrm{sub}}$
    for the saturated pair; 30~s, i.e.\ 22~$\tau_{\mathrm{sub}}$, for
    Molaro $-20\,^{\circ}$C); say so.
  \item \tA Conductivity sampled every step for $\theta < 11$, then at every
    0.1\% drop in SSA.
  \end{itemize}
\Pg{Checks}{to write}{The results do not depend on domain size, time step or sampling.}
  \begin{itemize}
  \item \tA One sentence each, pointing to the supplement: domain size
    (Fig.~S1); interface parameter (Fig.~S2), including that the first day
    depends on $\varepsilon$ and later rates do not; time step and sampling
    (Fig.~S3).
  \end{itemize}
\Place{Table 3}{\Have Complete parameters, placed. \tA Not cited in the text.
  \Open{keep it here, or the short version here and this one in the
  supplement.}}
\end{enumerate}

\subsection{Simulation Design: Grain Pairs and Aggregates}
\begin{enumerate}
\Pg{Grain pairs}{to write; heading only at L277; material now at L301--313 and L326--331}{What the three two-grain simulations are.}
  \begin{itemize}
  \item \tA Axisymmetric about the line of centers; fixed temperature and
    humidity on the outer boundary; how the neck width $w$ is measured.
  \item \tA The two Molaro geometries and the equal-grain pair ($R =
    84.4~\mu$m, $-20\,^{\circ}$C, 100~h).
  \item \tA Every pair starts from two overlapping spheres with a neck
    28~$\mu$m wide, not from grains that just touch. Say so here; it
    matters for the exponents (Section~3.2).
  \end{itemize}
\Place{Table 4}{\Have Two-grain parameters. \tA Add a column for the saturated
  pair. \tC Simulated time at $-5\,^{\circ}$C is 68~min, not 2~h.}
\Pg{Aggregates}{to write; material now at L343--352}{What the 125 aggregate simulations are and how the packings were made.}
  \begin{itemize}
  \item \tA 2~mm square, periodic in both directions; circular grains,
    log-normal radii with mean 50~$\mu$m, placed by drop-and-roll
    deposition.
  \item \tA Five porosities, five independent packings each, the same
    packings at all five temperatures: 125 runs of 30 days; uniform
    temperature, vapor initially at saturation.
  \item \tA \Chk{F5} Domain size $L/\bar R = 40$, \del{at the 1--2~mm representative volume
    of Calonne et al.\ (2011).} \add{from our domain-size test (F5); Calonne et al.\ (2011) give 2.5--5.5~mm for natural snow.} Packings are accepted only if contacts
    across the periodic seam are as dense as in the interior.
  \end{itemize}
  \Cite{representative volume: Calonne et al.\ 2011 \texttt{calonneNumericalExperimentalInvestigations2011}.}
\Pg{Why periodic, why two dimensions}{to write}{Justifies the two choices a reviewer will question first.}
  \begin{itemize}
  \item \tA Periodic: the domain stands for the interior of a deposit, no
    wall physics has to be invented, and the system is closed, so
    curvature alone drives transport.
  \item \tA Two dimensions: cost (now at L290--291, in the wrong place);
    what it costs is discussed in Section~5.6.
  \end{itemize}
\end{enumerate}

\section{Neck Growth Between Two Ice Grains}
\begin{enumerate}
\Pg{Section opening}{L279--294}{Announces the findings.}
  \begin{itemize}
  \item \tC It was written as the opening of a single Results section and
    still announces the aggregate findings of Section~4.
  \item \tC ``independent of the mean temperature as well as the initial
    porosity''; ``our model's ability to capture sintering dynamics for
    varying temperatures'': neither is what the results show.
  \item \tC ``theorized trends from Demmenie et al.'': they measured the
    exponent as well. ``dominated by vapor transport'': the $t^{1/3}$ law
    is for sublimation and desublimation limited by attachment.
  \item \tG ``simulation results, which captures'' (capture); ``related the
    evolution'' (to the); ``microstructures SSA'' (microstructure's).
  \item \tT Two sentences: before using the model on aggregates, we ask how
    much of measured neck growth it reproduces, with one experiment away
    from saturation (the rate) and one at saturation (the exponent).
  \end{itemize}
\end{enumerate}

\subsection{Comparison With Unsaturated Experiments}
\begin{enumerate}
\Pg{Why these experiments}{L295--300}{Why Molaro et al.\ is the comparison.}
  \begin{itemize}
  \item \tC ``relatively well-controlled'': the humidity was not reported
    and the grains sublimated; Section~5.1 argues the conditions were not
    controlled at the level of the driving force. Say ``among the few that
    isolate a grain pair at these sizes and temperatures''.
  \item \tG ``are good candidate'' (a good candidate); ``validating''
    (testing); ``that exist that isolates'' (isolate).
  \end{itemize}
\Pg{What is fitted and what is not}{L301--313}{The set-up of the $-20\,^{\circ}$C run.}
  \begin{itemize}
  \item \tC ``The only free-parameter in our model is \ldots\
    $\beta_{\mathrm{sub}}$'': the humidity is fitted too.
  \item \tC ``over a 78 minute period'': Molaro's text describes their
    Figure~11 as 57~min apart; our digitized series runs to 78~min. Check
    which figure the series comes from and cite it.
  \item \tC ``within a sealed cryostage'': their text says the sample area
    was not sealed off from the chamber, and the grains shrank.
  \item \tG ``fig.~11'' (their Figure~11); ``We assumed a value \ldots\ and
    was at the upper end'' (state $\alpha_c = 0.1$); ``(cite other works as
    well)''; ``table table 4''.
  \item \tA The inferred humidity, 0.99715 at $-20\,^{\circ}$C and 0.99674 at
    $-5\,^{\circ}$C, from the shrinkage of the large grain.
  \item \tT Geometry and boundary conditions move to Section~2.5.
  \end{itemize}
\Pg{Result}{L314--324}{Figure~3: model against measurement.}
  \begin{itemize}
  \item \tG ``Panel (a) we see'' (shows); ``the a digital replicate'';
    ``corresponds final temporal data point''; ``over approximates and
    under approximates'' (overpredicts, underpredicts); ``will discussed''
    (be); ``section section 5.1''; ``that they cryostage chamber''.
  \item \tA The numbers: 41\% of the measured neck growth at
    $-20\,^{\circ}$C over 78~min; 170\% at $-5\,^{\circ}$C over 48~min (52.1
    against 44.0~$\mu$m).
  \item \tA Panels (c), (d): the shrinkage is what the humidity was fitted
    to, so they are a calibration, not a test.
  \item \tA Molaro's conventions are used (neck width; time from the first
    measurement). Their own model matches through its vapor term, which
    assumes $\alpha_c = 1$ and no gas diffusion.
  \end{itemize}
  \Cite{Molaro et al.\ 2019 \texttt{molaroMicrostructuralEvolutionWater2019}.}
\Place{Figure 3}{\Have placed. Caption: see ``Figure captions'' after Section~6.}
\end{enumerate}

\subsection{Growth at Saturation}
\begin{enumerate}
\Pg{Set-up and result}{tex L454--456; \textbf{reframed 10-15} for the tangent-start run}{Two grains brought into contact in saturated air: the model gives the measured exponent.}
  \begin{itemize}
  \item \tA \textbf{Framing.} Section~3.2 is now the positive test of the
    two-grain comparison. At saturation the grains neither grow nor shrink
    (radius $+0.08\%$ over 150~h), so the neck is the only thing that grows,
    and the model gives the exponent Demmenie et al.\ measured and the
    sharp-interface law predicts. Section~3.1 (Molaro et al.) remains the
    test of the \emph{rate}; Section~3.2 is the test of the \emph{exponent}.
  \item \tA \textbf{Set-up, as run.} Two grains, $R = 84.4~\mu$m, just
    touching at $t = 0$; $-20\,^{\circ}$C; $\alpha_c = 1.002\times10^{-3}$;
    vapor saturated over the grains' own curvature, $\rho_v =
    \rho_{vs}^{I}(1 + 2d_0/R)$; 3-D axisymmetric; 150~h.
  \item \tA \textbf{Result.} Fitting $w = C(t+t_0)^{a}$ after the relaxation
    period (20~min, $11\,\tau_{\mathrm{sub}}$) gives $a = 0.289$ ($t_0 =
    0.6$~h, rms 0.17\%); the Kuczynski form gives $1/m = 0.290$; fixing
    $a = 1/3$ raises the rms to 2.6\%. Starting the fit later: 0.289 (1~h),
    0.285 (13~h), 0.281 (32~h), 0.279 (63~h), 0.277 (88~h). The neck reaches
    61.2~$\mu$m ($w/2R = 0.36$).
  \item \tA \textbf{What it is compared with.} Demmenie et al.\ measured
    0.26--0.33 at saturation (p.3), on spheres at $-3\,^{\circ}$C: the
    exponent is compared, not the rate. The sharp-interface law integrated
    for this geometry and $\alpha_c$ gives 0.29 (Appendix~B); $1/3$ is its
    limit for small necks, and the neck here is not small.
  \item \tA \textbf{Why a small $\alpha_c$ gives $t^{1/3}$, and a large one
    $t^{1/5}$} (the reader will ask; derivation in Appendix~B). Both are
    vapor transport; they differ in which step limits it. Vapor leaves the
    convex grain surface and arrives at the concave neck, driven by the
    curvature of the neck fillet, $\Delta\sigma \approx d_0/\rho$, with fillet
    radius $\rho \approx x^2/2R$. It passes two resistances in series:
    attachment at the ice surface, $\beta_{\mathrm{sub}} \propto 1/\alpha_c$,
    which does not depend on size, and diffusion through the air over a
    length of order $\rho$, which grows with the neck.
    \begin{itemize}
    \item Attachment-limited (small $\alpha_c$): $v_n \approx
      d_0/(\beta_{\mathrm{sub}}\rho) \propto R/x^2$, so $x^2\,dx \propto dt$,
      $x^3 \propto t$, $w \propto t^{1/3}$.
    \item Diffusion-limited (large $\alpha_c$): $v_n \propto
      d_0/(k_D\rho\cdot\rho) \propto R^2/x^4$, so $x^5 \propto t$, $w \propto
      t^{1/5}$.
    \end{itemize}
    A resistance that grows as the neck grows slows the neck faster, hence
    the smaller exponent. The classical $1/3$ (Kingery 1960, Table~I:
    ``Evaporation condensation, $(x/r)^3 = B(T)t/r^2$'') assumes the gas
    imposes no resistance: Kuczynski's model ``assumed the mean free path
    \ldots\ of water molecules at atmospheric pressure in air to be large
    compared with the neck radius'' (Blackford 2007, p.16). In air it is not
    ($\sim$0.1~$\mu$m), so Hobbs \& Mason put in diffusion through air and got
    $n = 5$ (same page). A small $\alpha_c$ makes attachment so slow that the
    gas resistance becomes negligible even in air, which restores the
    condition the $1/3$ law assumes. Here gas diffusion is 0.4--2\% of the
    resistance at $\alpha_c = 10^{-3}$; our neck ODE gives $a \approx 0.20$
    at $\alpha_c = 0.1$ and 0.29--0.30 at $10^{-3}$
    (\texttt{studies/molaro\_2019/envelope\_ode/}), and the run gives 0.289.
    It is below $1/3$ because the neck is not small ($x/R$ up to 0.36: the
    $2/R$ term and the exact fillet slow it).
  \item \tA \textbf{Grain scale against neck scale} (so this does not
    contradict Intro \P6). Bouvet et al.\ (2024, p.5) place the transition
    near $\alpha_c = 10^{-3}$ for 0.1~mm grains, with the diffusion length
    set by the grain size. At the neck the diffusion length is the fillet
    radius, a few $\mu$m, so the neck stays attachment-limited up to larger
    $\alpha_c$ than whole-grain growth does. In the aggregates the exchange
    is between grains, and the two resistances are comparable ($\Lambda
    \approx 2.8$).
  \item \tA \textbf{Where it goes:} the manuscript text is two sentences in
    Discussion~5.1 (``The saturated pair, and exponents''). This bullet and
    Appendix~B are the background for it, not text for the paper.
  \item \tC \Chk{F2} ``a vapor-limited regime \ldots\ governed by vapor
    transport'' (tex L455): the reverse; Bouvet et al.\ (2024, p.5) put the
    attachment--diffusion transition near $10^{-3}$ for 0.1~mm grains. At
    $\alpha_c = 10^{-3}$ growth is limited by attachment (gas diffusion adds
    0.4--2\% to the resistance), which is the regime the $t^{1/3}$ law
    describes.
  \item \tC Numbers in the draft that describe the earlier run: $t = 100$~h
    $\to$ 150~h; fit window $[20~\mathrm{min}, 100~\mathrm{h}]$ $\to$ $[20~\mathrm{min},
    150~\mathrm{h}]$; $a = 0.21$ $\to$ 0.29; ``briefly brought into
    contact'' $\to$ ``just touching''.
  \item \tC ``The red dashed curve is the theoretical fit'': it is the best
    fit with the exponent fixed at $1/3$.
  \item \tG ``in attempts to compare'' (to compare); ``allowed for the to
    sinter''; ``an condensation''; ``$1.002e-03$'' ($1.002\times10^{-3}$);
    ``Similarly to''.
  \Snippet{We performed one further simulation to compare the exponent of
    neck growth with that measured by
    \citet{demmenieIceSinteringSublimation2025}. Two ice grains of radius
    $R = 84.4~\mu$m, initially just touching, sinter at $T=-20^\circ$C for
    150 hours in vapor saturated over the grains' own curvature, so that the
    grains neither grow nor shrink and only the neck evolves. We use
    $\alphaC = 1.002\times10^{-3}$, for which growth is limited by attachment
    at the ice surface. After the initial relaxation period, the neck width
    follows $w = C(t+t_0)^{a}$ with $a = 0.29$
    (\cref{fig:Figure4_saturated_neck_growth}), and the exponent changes by
    less than 0.02 as the start of the fit is moved from 1 to 88 hours. This
    lies within the range $0.26$--$0.33$ that
    \citet{demmenieIceSinteringSublimation2025} measured at saturation, and
    matches the value 0.29 that the sharp-interface law gives for this
    geometry (Appendix~B); the small-neck limit of that law, $a = 1/3$, fits
    our data less well.}
  \end{itemize}
  \Cite{Demmenie et al.\ 2025 \texttt{demmenieIceSinteringSublimation2025}.}
\Pg{The starting neck}{to write; one or two sentences here, the rest in Appendix~B}{Why the run starts from touching grains.}
  \begin{itemize}
  \item \tA An earlier run started from a pre-made neck (two overlapping
    spheres, a 28~$\mu$m neck) and gave $a = 0.21$--0.24 under the same
    conditions. The low exponent came from that initial condition: a fit
    must assume how long a pre-made neck had been growing, and the free fit
    picks 4.8~h where the prediction says 16~h. Report it in Appendix~B as
    a caution on initial conditions, not in the main text.
  \end{itemize}
  \Open{the Molaro runs (Section~3.1) still start from a pre-made neck,
  because a tangent start traps an air pocket at the neck at that
  resolution (pilot of 10-15). Decide how Sections~3.1 and 5.1 qualify the
  0.135 exponent at $-20\,^{\circ}$C, which may carry the same bias, and
  whether the 41\% rate comparison, anchored on Molaro's first measurement,
  is affected.}
\Place{Figure 4}{\Have placed. Caption: see ``Figure captions'' after Section~6.}
\end{enumerate}

\section{Thermal Conductivity of Sintering Aggregates}

\subsection{Microstructural Evolution}
\begin{enumerate}
\Pg{Set-up and the gallery}{L343--361}{States the simulation matrix, then describes Figure~5.}
  \begin{itemize}
  \item \tC ``mesh of size $2829 \times 2928$'': 2829 both ways.
  \item \tC ``We initially consider a single 2D microstructure'': the
    previous sentence states 125 runs, and Figure~5 shows five.
  \item \tC ``periodic boundary conditions for our equation (1b) and (1c)'':
    all three fields are periodic.
  \item \tC ``By the 10 day snapshot for each porosity \ldots\ a well
    connected icy porous matrix'': the ice already spans the cell at $t =
    0$ in 24 of 25 packings; what changes is that point contacts become
    necks. At the two highest porosities the network stays anisotropic.
  \item \tG ``Next, we aim to better understand \ldots\ to better
    understand'' (twice); ``five different configurations for five
    separate porosities'' (five packings at each of five porosities);
    ``drop and roll'' (hyphens); ``gravitational affects'' (effects);
    ``a a''.
  \item \tT The set-up (first half) moves to Section~2.5.
  \item \tA What the colors show: vapor is undersaturated over convex ice,
    which sublimates, and supersaturated in necks and concavities, where
    it desublimates.
  \end{itemize}
\Place{Figure 5}{\Have placed. Caption: see ``Figure captions'' after Section~6.}
\end{enumerate}

\subsection{Temperature Rescales Time}
\begin{enumerate}
\setcounter{enumi}{1}
\Pg{Computing the conductivity}{L362--371}{Applies the homogenization and describes the tensor.}
  \begin{itemize}
  \item \tC ``We obtain \ldots\ $\mathbf{K}_{\mathrm{eff}}$ by integrating
    (17a) and applying the zero mean condition'': it is the volume average
    of $\mathbf{K}(\nabla\mathbf{t} + \mathbf{I})$ over the cell (Eq.~10 in
    diffuse form).
  \item \tC ``off diagonal components are at least one order of magnitude
    lower'': $k_{xy}$ reaches about 12\% of $k_{xx}$ on one packing; it is
    zero in the mean.
  \item \tG ``changes \ldots\ is expected'' (are); ``figure figure 6'';
    ``correspond the the''; ``lower that'' (than); ``over various seeds''
    (packings).
  \item \tT The two sentences on solving Eq.~(17) on the same mesh repeat
    Section~2.3.
  \item \tA Anisotropy: $k_{xx}/k_{yy} = 1.02$, 1.03, 0.93, 0.82, 0.66 from
    the lowest porosity to the highest, with $y$ the direction of
    deposition. The scalar $k_{\mathrm{eff}}$ is used from here on.
  \end{itemize}
\Pg{Figure 6, one packing at five temperatures}{L372--381}{Conductivity and SSA against time and against sintering age.}
  \begin{itemize}
  \item \tC ``After the first day'': that is 11~$\tau_{\mathrm{sub}}$ at
    $-20\,^{\circ}$C only; the relaxation period is a fixed number of
    $\tau_{\mathrm{sub}}$, so it is 0.26~d at $-5\,^{\circ}$C and 7.7~d at
    $-40\,^{\circ}$C.
  \item \tC ``the trend is linear'': linear in $\log\theta$, about 0.13 per
    decade.
  \item \tG ``local sintering occurs at high rates \ldots\ and the
    phase-field interface relaxes'' joins two different things;
    ``non-dimensional'' here, ``nondimensional'' in the caption.
  \item \tA Panel (a) is not mentioned.
  \item \tA The result as a number: against $\theta$ the rates follow
    $\tau_{\mathrm{sub}}$ within 4\%; against SSA the five temperatures
    agree within 0.27\% for every one of the 25 packings (median 0.07\%).
  \item \tA A size that does not depend on $\varepsilon$: between $\theta =
    30$ and 331, $k_{\mathrm{eff}}$ rises by a factor $1.171 \pm 0.015$ and
    SSA falls to $0.819 \pm 0.007$, for $\varphi \le 0.375$.
  \item \tA Why curves are normalized at $\theta_r = 30$ and not at $t = 0$:
    the first day depends on the interface parameter (Fig.~S2).
  \end{itemize}
  \Cite{for SSA against log time in isothermal snow: Legagneux \& Domin\'e 2005 \texttt{legagneuxMeanFieldModel2005}; Taillandier et al.\ 2007 \texttt{taillandierRateDecreaseSpecific2007} (functional form only; verify).}
\Place{Figure 6}{\Have placed. Caption: see ``Figure captions'' after Section~6.}
\end{enumerate}

\subsection{Conductivity as a Function of Specific Surface Area}
\begin{enumerate}
\setcounter{enumi}{3}
\Pg{Figure 7, porosity}{L382--400}{One paragraph covering all four panels.}
  \begin{itemize}
  \item \tC ``15 unique packings--3 packings for each porosity'': five
    packings at each of three porosities.
  \item \tC ``We show the time evolution of $k_{\mathrm{eff}}$ for each of
    the simulations'': panel (a) is $-20\,^{\circ}$C only, the mean of five
    packings and their range.
  \item \tC ``the ice matrix remains disconnected in the horizontal
    direction'': not so. The ice spans the cell in both directions in every
    packing from 10~d on (percolation test). What happens is that the
    packings become anisotropic and the two directions follow different
    laws: along the deposition direction the exponent stays at $-0.79$ to
    $-0.86$ up to $\varphi = 0.425$; across it, $-0.70$, $-0.60$, $-0.48$
    from $\varphi = 0.375$ on.
  \item \tC ``The linear trend'' in panel (c): the axis is logarithmic; the
    fit is exponential in porosity.
  \item \tG ``may effect'' (affect); ``scientifically their connectivity
    as.'' is unfinished; ``the strength of the reference value'' (the
    reference value); ``each of the the''; ``A single point \ldots\ show''
    (shows); ``panels(c)''.
  \item \tT One paragraph carries four panels; split into three: (a); (b);
    (c) and (d).
  \item \tA Why four temperatures: the $-40\,^{\circ}$C runs end at $\theta =
    43$, too soon after the reference age.
  \item \tA The numbers. Exponent $-0.80$; by porosity $-0.79$, $-0.82$,
    $-0.78$, then $-0.73$, $-0.63$. Level $k_{\mathrm{eff},r}/
    k_{\mathrm{ice}} = 1.37\exp(-5.01\,\varphi)$ for $\varphi \le 0.375$,
    6.6\% scatter between packings. At 30~d the conductivity differs by a
    factor of about 3.3 across porosity.
  \item \tA What it amounts to: SSA gives the relative change in
    conductivity; porosity gives the level; neither alone gives both.
  \end{itemize}
  \Open{whether to plot or tabulate the directional exponents; whether to
  add that a porosity-dependent representative volume may contribute (not
  tested: the domain-size study was run at $\varphi = 0.325$ only).}
\Place{Figure 7}{\Have placed. Caption: see ``Figure captions'' after Section~6.}
\Pg{What the scatter is}{optional; to write}{Differences between packings in the exponent are properties of the packings, not noise.}
  \begin{itemize}
  \item \tA The scatter in the exponent does not fall with domain size while
    the scatter in the level does (11\% to 4.5\%).
  \end{itemize}
  \Open{include here, in Section~5.3, or leave out.}
\end{enumerate}

\section{Discussion}
Paragraphs 1--3 of Section~5.1 and the start of Section~5.2 are drafted.
\begin{enumerate}[label=(\alph*)]
\Pg{Section opening}{L402--411}{Recaps the paper.}
  \begin{itemize}
  \item \tC ``our model is able to capture the trends in sintering dynamics
    at multiple temperatures'': at $-20\,^{\circ}$C it gives 41\% and at
    $-5\,^{\circ}$C 170\%.
  \item \tT The third statement of the plan (after the Introduction and the
    opening of Section~3). Not needed.
  \end{itemize}
\Pg{Lineage of the model}{L412--424}{Kaempfer \& Plapp, Moure \& Fu, and what this paper adds.}
  \begin{itemize}
  \item \tG ``They also, qualitatively validate''; ``has since become
    popular method''.
  \item \tT Belongs in the Methods opening, in one or two sentences.
  \end{itemize}
\end{enumerate}

\subsection{What Vapor Transport Alone Reproduces}
\begin{enumerate}
\Pg{The unsaturated pair}{L426--452}{Why the model gives 41\% of Molaro's neck growth and what that does and does not mean.}
  \begin{itemize}
  \item \tC ``underpredicts \ldots\ overpredicts \ldots\ by 41\% and 170\%'':
    the model gives 41\% \emph{of} the measured growth at $-20\,^{\circ}$C
    (it underpredicts by 59\%) and 170\% of it at $-5\,^{\circ}$C.
  \item \tC ``It does not account for surface diffusion, which is expected
    [to] contribute significantly'' (twice): by Molaro's own equations
    surface diffusion is below 1\% of vapor transport at these neck sizes
    (Appendix~A). Remove. Molaro et al.\ say so themselves (p.30: ``Grains
    100~$\mu$m in radius showed no significant effects from neglecting
    surface diffusion on any timescales''); their grains are 72--101~$\mu$m
    in radius. \Chk{F3}
  \Snippet{At $T=-20\,^\circ$C our model gives 41\% of the measured neck
    growth, and at $T=-5\,^\circ$C it gives 170\%. Our model contains no
    surface diffusion, but by the equations and parameters of
    \citet{molaroMicrostructuralEvolutionWater2019}, surface diffusion adds
    less than 1\% to vapor transport at these neck sizes
    (\cref{sec:Appendix_SD_VT}), so it does not explain the shortfall at
    $-20\,^\circ$C.}
  \item \tC ``This result was expected'': it was not.
  \item \tC ``The chamber being unsaturated \ldots\ [is] likely the reason
    why our model predictions overestimate'': the undersaturation is
    already in the model, through the inferred humidity.
  \item \tG ``preformed'' (performed); ``the the''; ``our model own
    accounts'' (only accounts); ``expected contribute''; the three
    sentences from ``Furthermore, our model also underestimates'' appear
    twice.
  \item \tT The first three sentences repeat Section~3.1.
  \item \tA ``Even with the vapor diffusivity raised 30 or 100 times'': say
    these were tests not shown, and that the neck was insensitive to the
    chamber humidity at the nominal diffusivity.
  \item \tA How small the driving force is: $4.8\times10^{-4}$ falling to
    $9\times10^{-5}$ in relative vapor density, the equivalent of 4.9 to
    1.0~mK; the chamber's undersaturation is $2.85\times10^{-3}$, 6 to 31
    times larger.
  \item \tA What that makes the result sensitive to, without choosing:
    (i)~how much of the ambient undersaturation reaches the neck (about 6\%
    in our run; with none, the one-equation neck model gives 66\% of the
    growth and an exponent of 0.23, the value of our fit to Molaro's data); (ii)~temperature
    differences of millikelvin between neck and grain surface;
    (iii)~$\alpha_c$ above 0.1 for droplets quenched in liquid nitrogen.
  \item \tA Latent heat is not a candidate: in the run the neck is 0.002~mK
    warmer than the grain surface.
  \item \tA Molaro's own model matches through a vapor term with $\alpha_c =
    1$ and no gas diffusion, and they present its agreement as
    order-of-magnitude.
  \item \tA Why the saturated experiment is the better test: it removes the
    undersaturation.
  \end{itemize}
  \Cite{Molaro et al.\ 2019 \texttt{molaroMicrostructuralEvolutionWater2019}; Swinkels \& Ashby 1981 \texttt{swinkelsSecondReportSintering1981}.}
\Pg{The saturated pair, and exponents}{tex L502; \textbf{reframed 10-15}}{The model reproduces how the exponent depends on saturation; the unsaturated value carries the starting-neck caveat.}
  \begin{itemize}
  \item \tA \textbf{Framing.} The model reproduces the dependence of the
    exponent on saturation that Demmenie et al.\ observed: 0.29 at
    saturation (theirs 0.26--0.33) and 0.135 in the undersaturated
    $-20\,^{\circ}$C case (theirs about $1/7$ when the ice sublimated). The
    draft's ``they are lower than measured values'' no longer holds at
    saturation.
  \item \tC ``As two spherical ice grains sinter in an environment with a
    saturated water vapor field, the growth is expected to be kinetics
    limited'': saturation does not make growth attachment-limited; a small
    $\alpha_c$ does. The $t^{1/3}$ law assumes it.
  \item \tA \textbf{Why $1/3$ at small $\alpha_c$} (replaces the sentence
    above; reasoning in Section~3.2 and Appendix~B). Two sentences:
  \Snippet{An exponent near $1/3$ is expected only when attachment at the
    ice surface, rather than diffusion of vapor through the air, limits the
    growth of the neck: the resistance to attachment does not depend on the
    size of the neck, whereas the diffusion path lengthens with the radius of
    curvature of the neck, which lowers the exponent to $1/5$
    \citep{kingery_regelation_1960, Hobbs_and_Mason,
    blackfordSinteringMicrostructureIce2007}. With $\alphaC = 10^{-3}$,
    diffusion through the air contributes less than 2\% of the resistance at
    the neck, so our saturated case lies in the regime the $t^{1/3}$ law
    describes.}
  \item \tC Numbers: ``$a = 0.21$ across the entire simulation \ldots\ as
    high as $a = 0.236$'' describe the pre-necked run: $a = 0.29$, 0.277 to
    0.289 as the window moves.
  \item \tC \Chk{F4} ``a much better match for the theorized scaling \ldots\
    $t^{1/7}$'' (tex L502): $1/7$ is not a theory of unsaturated sintering.
    It is the classical surface-diffusion exponent (Kingery 1960, Table~I)
    and the value Demmenie et al.\ measured when ``we deliberately set the
    humidity too low and let the ice sublimate'' (p.3). Our model reaches it
    with vapor transport alone. Add that this run starts from a pre-made
    neck (Section~3.2, open item).
  \item \tA ``We additionally compare our measured interface velocity
    \ldots'' states no result. The numbers in Appendix~B (2.5 times the
    Gibbs--Thomson value at 2~h, 1.0 by 71~h) are for the pre-necked run;
    \textbf{recompute them for the tangent-start run}
    (\texttt{interface\_width\_estimate.py}) or drop the sentence.
  \item \tA Why not surface diffusion: mechanisms in parallel add their
    rates, and surface diffusion alone gives $t^{1/7}$, so it could not
    raise an exponent; at saturation none is needed.
  \item \tG The first sentence is unfinished (``growing under as .'');
    ``where observe''; ``we had considered''; ``$1.002e-03$''; ``preform''
    (perform); two sentences in a row begin ``We additionally'' and
    ``Additionally''.
  \Snippet{At saturation, the grains neither grow nor shrink and the neck
    width follows $w \propto t^{0.29}$ (Section~3.2), within the range that
    \citet{demmenieIceSinteringSublimation2025} measured, $0.26$--$0.33$. In
    the undersaturated case at $T=-20^\circ$C, where the grains sublimate,
    we measure $a = 0.135$, close to the exponent of about $1/7$ that
    \citet{demmenieIceSinteringSublimation2025} measured when the humidity
    was too low and the ice sublimated. The same value is the classical
    exponent for surface diffusion \citep{kingery_regelation_1960}, which our
    model does not contain; here it arises from vapor transport alone. Our
    model therefore reproduces how the exponent depends on the saturation of
    the vapor.}
  \end{itemize}
  \Cite{Demmenie et al.\ 2025 \texttt{demmenieIceSinteringSublimation2025}; Molaro et al.\ 2019 \texttt{molaroMicrostructuralEvolutionWater2019}; Kingery 1960 \texttt{kingery\_regelation\_1960}.}
\Pg{$\alpha_c$ as an effective parameter}{L469--477}{What the uncertain coefficient does and does not change.}
  \begin{itemize}
  \item \tC ``do not change how the microstructure evolves, just the rate'':
    true only while growth is limited by attachment. $\alpha_c$ sets the
    time scale ($\tau_{\mathrm{sub}} \propto 1/\alpha_c$) \emph{and} the
    group $\Lambda = L^{*}/\bar R \propto 1/\alpha_c$ that decides how much
    gas diffusion matters. In the aggregates $\Lambda \approx 2.8$. A
    smaller $\alpha_c$ only rescales time; a larger one brings diffusion
    in and changes the path as well. (An earlier version of this outline
    said a tenfold larger $\alpha_c$ would leave the curves unchanged
    against $\theta$; that was wrong.)
  \item \tC ``$\Delta\beta/\beta_{\mathrm{sub}} = 0.02$'' for both the
    aggregates and the saturated pair: it is 0.02, 0.002, and 0.23 for
    the Molaro pairs.
  \item \tC \Chk{F7} The sentence on Bouvet et al.\ rescaling time through
    $\alpha_c$ (tex L504) cites the 2023 paper, a laboratory experiment with
    no model. It is Bouvet et al.\ (2022, p.2: ``Relating the numerical time
    \ldots\ to the real physical time involves the condensation
    coefficient''). Their model is limited by attachment alone, so the
    rescaling is exact there; in ours it holds while attachment limits
    growth (previous bullet).
  \Snippet{\citet{bouvetSnowEquiTemperatureMetamorphism2022} used this to
    calibrate their model: matching the SSA evolution measured by $\mu$CT
    fixed $\alphaC$, and with it the physical time of their simulations. In
    their model growth is limited by attachment alone, so the rescaling is
    exact; in ours it holds while attachment, not vapor diffusion, limits
    growth.}
  \item \tG ``The kinetic attachment parameter, and thus condensation
    coefficient, do not''; ``$\mu$m-CT'' ($\mu$CT).
  \item \tT Eq.~($\tau_{\mathrm{sub}}$) and the sentence after it move to
    Section~2.2.
  \item \tA What $\alpha_c$ does to the two grain-pair comparisons: Molaro's
    rate wants a large value (41\% at 0.1; 51--75\% at 0.8 in the
    one-equation neck model), while the sharp-interface exponent at
    saturation falls with it (0.29 at $10^{-3}$, 0.25 at $10^{-2}$, 0.18
    at 0.1). One value cannot serve both.
  \item \tA \Note Bouvet et al.\ fit $\alpha_c = (9.8 \pm 0.7)\times10^{-4}$
    to tomography at $-2\,^{\circ}$C; Braun et al.\ find
    $10^{-3}$--$10^{-1}$, changing as surfaces facet, and call it an
    effective, surface-averaged value.
  \end{itemize}
  \Cite{Bouvet et al.\ 2022 \texttt{bouvetSnowEquiTemperatureMetamorphism2022}; Braun et al.\ 2024 \texttt{braunRigorousApproachSpecific2024}; Granger et al.\ 2021 \texttt{grangerOrientationSelectiveGrain2021} for faceting under a gradient; Libbrecht \& Rickerby 2013 \texttt{libbrechtMeasurements\allowbreak SurfaceAttachment2013} as background on why $\alpha_c$ is uncertain.}
\end{enumerate}

\subsection{Temperature as a Rescaling of Time}
\begin{enumerate}
\setcounter{enumi}{3}
\Pg{Why the curves superpose}{L479--487; breaks off at ``Since we have already''}{A similarity argument: in dimensionless form, almost all of the temperature dependence is in one time scale.}
  \begin{itemize}
  \item \tC ``$\beta_{\mathrm{sub}}$, which is a free variable'': it is set
    by $\alpha_c$, which is chosen, and it depends on temperature through
    $\rho_{vs}^{I}(T)$. That dependence is the point of the paragraph.
  \item \tG ``is evident \ldots\ and is evident''; ``a two part series''
    (two steps in series); ``drawn to areas of lower vapor saturation''.
  \item \tA Step 1. Scale lengths by $\bar R$, supersaturation by
    $d_0/\bar R$, time by $\tau_{\bar R} = \bar R^{2}\beta/d_0$
    ($= 2500\,\tau_{\mathrm{sub}}$). The interface condition becomes
    parameter-free. The saturation vapor density, with its 31-fold change
    from $-40$ to $-5\,^{\circ}$C, is entirely inside the time scale.
  \item \tA Step 2. One temperature-dependent group is left, $\Lambda =
    L^{*}/\bar R$ with $L^{*} = (D_v/\alpha_c)\sqrt{2\pi m/k_BT}$. The
    saturation vapor density cancels in it; what remains is the
    diffusivity. It changes by 20\% where the time scale changes by a
    factor of 29.
  \item \tA Step 3. So the saturation vapor density sets the clock, and the
    diffusivity is a weak remaining correction: the likely source of the
    4\% by which rates depart from $\tau_{\mathrm{sub}}$ (consistent in
    size; not tested as the cause).
  \item \tA Say plainly that with one sintering mechanism this is close to a
    consequence of the construction; the numerical result is how well it
    holds. Do not say gas diffusion is negligible ($\Lambda \approx 2.8$);
    what is small is its change with temperature.
  \item \tA The two-page derivation
    (\texttt{temperature\_scaling/\allowbreak time\_temperature\_scaling.pdf}) can be
    an appendix or supporting information.
  \end{itemize}
\item \New{What would break it, and the activation energy.} \Miss
  \Msg{The collapse is therefore a signature of a single route. Its temperature dependence differs from the one measured for strength gain, which points to a second route in real samples and gives experimenters something to measure.}
  The
  clock's temperature dependence is an activation energy of about
  50~kJ\,mol$^{-1}$, that of the saturation vapor density. \Note Choukroun
  et al.\ measure $24.3 \pm 3.3$~kJ\,mol$^{-1}$ for strength gain and
  attribute it to surface self-diffusion (L613--620), at 193--243~K; the two
  temperature ranges overlap at 233--243~K. The difference is consistent
  with a second route in their samples, and an activation energy is
  something a laboratory can measure to tell the routes apart. A temperature-dependent
  $\alpha_c$, a second route, or a temperature gradient would each break the
  superposition, which makes it a test.
  \Cite{Choukroun et al.\ 2020 \texttt{choukroun\_strength\_2020}.}
\end{enumerate}

\subsection{Specific Surface Area as the State Variable}
\begin{enumerate}
\setcounter{enumi}{5}
\item \New{Density is not enough.} \Miss
  \Msg{Conductivity changes while porosity stays fixed, so a density-based formula cannot describe a sintering deposit.}
  In these runs porosity is fixed and
  conductivity still rises, by 17\% between $\theta = 30$ and 331. \Note Calonne et al.\ (2011): tensor from
  30 tomographic images, a regression on density, scatter attributed to
  anisotropy (L507--517). Calonne et al.\ (2014): under a temperature
  gradient conductivity changes at constant density, with SSA and anisotropy
  (L517--529). Half the length of the notes.
  \Cite{Calonne et al.\ 2011 \texttt{calonneNumericalExperimentalInvestigations2011}; Calonne et al.\ 2014 \texttt{calonneStudyTemperatureGradient2014}; Schneebeli \& Sokratov 2004 \texttt{schneebeliTomographyTemperatureGradient2004} (the notes cite it for the same behavior under a gradient).}
\item \New{The power law and where it stops.} \Miss
  \Msg{State the law, its range, and that outside the range we make no claim.}
  $p = -0.80$ for
  $\varphi \le 0.375$; the level from porosity; above 0.375 the exponent
  weakens and the packings are anisotropic. No claim that packings of
  different porosity converge.
  \Cite{Calonne et al.\ 2011, 2019 \texttt{calonneNumericalExperimentalInvestigations2011}, \texttt{calonneThermalConductivitySnow2019} for conductivity against density, as the contrast.}
\item \New{Why packings differ.} \Miss
  \Msg{Packings of the same porosity conduct differently because their ice networks are arranged differently; at high porosity the networks are anisotropic and the simple law fails in one direction.}
  Packings of one porosity differ in
  how well the ice is connected, and the exponent follows that (correlation
  $-0.77$ with conduction efficiency), not coordination number or
  anisotropy. In these packings the ice always spans the cell and the pore
  never does (percolation test, 2026-10-12): the pore space is a set of
  closed pockets, which is what ``sealed pores'' means here, and in two
  dimensions only one phase can span. What differs between packings is the
  quality of the ice network, not whether it exists: at high porosity the
  packings are anisotropic, grains dropped under gravity being supported
  from below with fewer contacts sideways, and the power law survives along
  the deposition direction while failing across it. Absolute values are
  therefore not compared with three-dimensional measurements, where both
  phases are connected.
  \Open{literature thresholds for discs are not yet checked against a
  source; whether to show the efficiency; the Euler-characteristic series
  (tested on three runs) is available if a direct measure is wanted.}
\end{enumerate}

\subsection{What a Laboratory Could Measure}
\begin{enumerate}
\setcounter{enumi}{8}
\item \New{Three tests.} \Miss
  \Msg{Turn the results into measurements someone could make on sintered ice-sphere samples.}
  On unloaded ice-sphere samples held
  isothermally: (i)~curves at different temperatures superpose under a time
  shift, and the shift gives an activation energy; (ii)~conductivity against
  SSA is one curve; (iii)~SSA falls linearly in log time. \Note SSA is
  measured by near-infrared reflectance, gas adsorption and tomography
  (L562--568; Gallet, Matzl, Proksch).
  \Cite{Gallet et al.\ 2009 \texttt{galletMeasurementSpecificSurface2009}; Matzl \& Schneebeli 2006 \texttt{matzlMeasuringSpecificSurface2006}; Proksch et al.\ 2015 \texttt{prokschDensitySpecificSurface2015}; for isothermal SSA decay in the laboratory, Legagneux \& Domin\'e 2005, Taillandier et al.\ 2007, Kaempfer \& Schneebeli 2007 \texttt{kaempferObservationIsothermalMetamorphism2007}, L\"owe et al.\ 2011 \texttt{loweInterfacialStructuralRelaxations2011} (forms only).}
\end{enumerate}

\subsection{Implications for Plume Deposits on Enceladus}
\begin{enumerate}
\setcounter{enumi}{9}
\item \New{What carries over.} \Miss
  \Msg{For Enceladus the qualitative conclusion holds: a sintering deposit conducts better as its surface area falls, even with no compaction.}
  Qualitative: as a deposit sinters, SSA
  falls and conductivity rises together, without densification. \Note Tiger
  Stripe temperatures up to about 180~K (Spencer); grain size increasing
  toward the fractures in VIMS data (Jaumann; L580--602); Molaro et al.\ on
  the dependence of sintering on temperature and grain size.
  \Cite{Spencer et al.\ 2009 \texttt{spencerEnceladusActiveCryovolcanic2009}; Spencer \& Nimmo 2013 \texttt{spencerEnceladusActiveIce2013}; Jaumann et al.\ 2008 \texttt{jaumannDistribution\allowbreak IcyParticles2008}; Molaro et al.\ 2019; Choukroun et al.\ 2020.}
\item \New{Why it matters, and the gap.} \Miss
  \Msg{Conductivity of the surface layer affects thermal and radar interpretation, but our conditions are far from the moon's, so we give no timescales.}
  \Note A low-conductivity
  surface layer changes subsurface temperatures and radar attenuation
  (Byrne; L603--612); Choukroun et al.\ on strength gain and fallout rates
  of 0.01--1~mm\,yr$^{-1}$ (L613--629). Then the gap, stated plainly: the
  runs stop at $-40\,^{\circ}$C and use 50~$\mu$m grains in air; the moon is
  far colder, the grains are micron-sized, and there is no air. No
  timescales are given.
  \Cite{Byrne et al.\ 2024 \texttt{byrneRadarAttenuationSubsurface2024} (a preprint DOI in the \texttt{.bib}; check whether it is published); Choukroun et al.\ 2020; possibly Ferrari et al.\ 2021 \texttt{ferrariThermalEmissionSaturns2021} on thermal emission of Saturn's icy moons (in your \texttt{.bib}; check relevance).}
\end{enumerate}

\subsection{Limitations}
\begin{enumerate}
\setcounter{enumi}{11}
\item \New{Model.} \Miss
  \Msg{What physics is missing.}
  Vapor route only; constant $\alpha_c$; air-filled,
  sealed pores at 1~atm; conduction only and constant $k_{\mathrm{ice}}$,
  $k_{\mathrm{air}}$. \Note Fourteau et al.\ on the contribution of vapor
  transport to apparent conductivity under a gradient (L569--575; check the
  figure quoted); Calonne et al.\ (2019) on temperature-dependent
  conductivities (L575--579).
  \Cite{Fourteau et al.\ 2021 --- your \texttt{.bib} has two, \texttt{fourteauImpactWaterVapor2021} and \texttt{fourteauMacroscopicWaterVapor2021}; the point about apparent conductivity is the first. Calonne et al.\ 2019 \texttt{calonneThermalConductivitySnow2019}.}
\item \New{Numerics and geometry.} \Miss
  \Msg{What the simulations cannot resolve or represent.}
  Two dimensions, so relative
  changes are claimed and absolute conductivities are not; one grain size
  distribution; rounded grains (discs) in point contact, so the rounding
  of fresh, dendritic snow, the fast first stage of isothermal
  metamorphism (Kaempfer \& Schneebeli 2007, p.1, p.5), is outside the
  model; the first day depends on the interface parameter; the Molaro
  runs start from a neck put in by hand, which biased the fitted exponent
  low in the saturated case (0.21 against 0.29 from contact; Appendix~B); the Molaro runs use time steps of 22 and 49
  $\tau_{\mathrm{sub}}$, outside the range the time-step study covered; pore
  throats narrower than the diffuse band (9.2~$\mu$m) are not resolved,
  which raises $k_{\mathrm{eff}}$ most at low porosity; the power law is
  established for $\varphi \le 0.375$; rates scale with $1/\alpha_c$, whose
  value is uncertain over two decades.
  \Cite{none needed beyond those above.}
\end{enumerate}

\textbf{In the notes and not to be carried over:} $\beta_{\mathrm{sub}}(T)$
from Libbrecht's $\sigma_0(T)$, a single tuning and a ``genuine prediction''
at $-5\,^{\circ}$C (L454--481); isotropy, and a collapse across porosity
(L530--561); sintering ``on the order of decades'' (L596) and conductivity
starting ``near 0.0'' (L611).

\section{Conclusions}
\begin{enumerate}
\item \New{What was done and found.} \Miss
  \Msg{Three results, each with the range over which it was shown.}
  One sentence of method, then the
  three results with their ranges: temperature rescales time ($-40$ to
  $-5\,^{\circ}$C, within 0.3\%); conductivity follows SSA$^{-0.80}$ with the
  level set by porosity ($\varphi \le 0.375$); vapor transport gives 41\% of
  measured neck growth at $-20\,^{\circ}$C.
\item \New{What comes next.} \Miss
  \Msg{The experiment and the model extensions that would test or extend them.}
  The laboratory test of Section~5.4; a
  second transport route in the model; three dimensions; colder
  temperatures and smaller grains.
\end{enumerate}

\section*{Figure captions}
Checked against the figure files in the 10-14 upload (rendered here), not
only against the text. AGU asks that a caption let the figure be read on its
own: what is shown, the conditions, what each line, symbol and color means.
Four points apply to several figures:
\begin{itemize}
\item \tC \textbf{The $\sigma$ color scale} (Figures 3--6) is relative to a
  flat ice surface and is clipped. Say both, and give the range. Without this
  a reader cannot tell why the pore is dark red at saturation (Figure~4) or
  uniformly dark blue in Molaro's chamber (Figure~3). The scale is
  symmetric-logarithmic (ticks at $\pm0.01$ and $\pm$ the maximum): say so.
\item \tG \textbf{``vapor saturation filed is colored according the the
  blue/red colorbar''} (Figures 3, 4): ``the supersaturation $\sigma$ is
  shown on the blue--red scale''. ``pale-blue'' and ``pale blue'' both occur:
  one form.
\item \tG \textbf{The phase-field colorbar is labeled $\phi_i$} in the art
  of Figures 3, 4, 6 and in the Figure~2 caption; the text writes $\phi$.
  One symbol (decided 09-28: $\phi$).
\item \tA \textbf{Which run.} Every image caption should name the porosity,
  temperature and domain of what is shown; most do not.
\end{itemize}
\begin{enumerate}
\Pg{Figure 1, mechanisms and an aggregate}{L268}{(a) a grain pair before and after necking; (b) an aggregate at four times.}
  \begin{itemize}
  \item \tC \textbf{The art says $R(t)$, the caption $w(t)$.} The arrow runs
    from the axis to the neck surface, i.e.\ a neck \emph{radius}, while $w$
    is the neck \emph{width} everywhere else (Figures 3, 4; Molaro's
    convention). Relabel the arrow $w(t)$ and draw it across the whole neck,
    or label it $w/2$.
  \item \Have The caption is vapor-only (10-15). \tG ``Ice growth as a
    result fo vapor transport is shown'': ``of''; or ``Arrows show vapor
    transport''. It still says ``a neck of radius $w(t)$'' (next bullet).
  \item \tA (b): which run (porosity, temperature, seed) and that the panels
    are at 0, 10, 19 and 30 days; the scale bar is 100~$\mu$m (it is a
    crop, so say ``a region of'').
  \item \tG ``A 2-D cross-section of two spherical ice grains that have
    undergone sintering. The initially non-sintered touching ice grains
    (shown as blue dotted line)'': ``Two ice grains in cross-section, before
    (dotted) and after (solid) a neck of width $w(t)$ has grown''. The grey
    crosshair lines in (a) are unexplained: remove, or say they are the axes
    of symmetry.
  \end{itemize}
\Pg{Figure 2, the homogenization}{L(after 2.3)}{(a) a cell, (b) the interface and mesh, (c) the corrector, (d) the heat flux.}
  \begin{itemize}
  \item \tA Not cited in the text (Section~2.3 is the place).
  \item \tA Which sample: porosity, temperature, age, and that the cell is
    2~mm wide (the scale bar is 500~$\mu$m).
  \item \tC (b): the dashed lines are orange, not red, and are the
    $\phi = 0.05$ and $0.95$ contours; the solid line is $\phi = 0.5$.
    Write $\phi$, not $\phi_i$.
  \item \tA (c): the black contours are the ice--air interface; only the
    $x$ cell problem is shown ($\mathbf{t}$ has a second component from the
    $y$ problem). Missing period after (c).
  \item \tA (d): color is on a logarithmic scale; ``per unit mean gradient
    in the $x$ direction'' is right; ``brightness of the path'':
    ``brighter paths carry more heat''.
  \item \tG ``the mesh in which we solve our homogenization model over'':
    ``the mesh on which the cell problem is solved''.
  \end{itemize}
\Pg{Figure 3, Molaro et al.}{L(Section 3.1)}{(a) the pair at the start and at 78 min, (b) neck width, (c), (d) grain diameters.}
  \begin{itemize}
  \item \tC \textbf{The color scale does not show the chamber.} The far
    field is undersaturated by $\sigma \approx -2.9\times10^{-3}$ (99.715\%
    relative to flat ice), a hundred times beyond the $\pm0.28\times10^{-4}$
    scale, so the whole pore sits at the bottom color. Say the scale is
    clipped to resolve the neck, and give the far-field value.
  \item \tC ``On the left is the initial configuration'': it is the first
    measured frame ($t = 0$ is Molaro's first measurement, and the
    simulation starts from the measured neck). Say so; it is why the curves
    in (b) start on the first data point. ``78 minutes'': check against
    Molaro's text, which describes their Figure~11 as 57 minutes apart
    (Section~3.1 entry).
  \item \tA (b): both temperatures are shown ($-20\,^{\circ}$C blue,
    $-5\,^{\circ}$C vermilion); lines are the model, circles the data, error
    bars are Molaro's stated uncertainty (check what they are); the numbered
    circles mark the frames in (a). Give the grain diameters.
  \item \tC (c), (d): ``the experimentally measured grain size through time
    relative to the large and small grains'': they show the diameter of the
    large (c) and small (d) grain divided by its value at $t = 0$; solid
    lines are the model, circles the data, dashed lines linear fits to the
    data.
  \item \tT ``The simulation is designed to replicate the conditions and
    geometry \ldots'': belongs in Section~3.1; the caption only needs the
    conditions.
  \end{itemize}
\Pg{Figure 4, the saturated pair}{L(Section 3.2); \textbf{rebuilt 10-15} from the tangent-start run}{(a) the pair at 0 and 150~h, (b) neck width with the fitted law ($a = 0.29$) and the $1/3$ law.}
  \begin{itemize}
  \item \tC The caption describes the old figure: 100~h $\to$ 150~h; the
    $t = 0$ frame now shows two grains just touching.
  \item \tA \textbf{Why the pore is red.} $\sigma$ is relative to flat ice,
    and Section~3.2 sets the far field to the Gibbs--Thomson value of the
    grains, $2d_0/R \approx 2.4\times10^{-5}$ for $R = 84~\mu$m, the top of
    the scale. Say this; a reader expects a saturated pore to be white.
  \item \tA The legend: $C(t+t_0)^{a}$, $a = 0.29$, and $C(t+t_0)^{1/3}$.
    Say both are fitted ($t_0$ free), over 20~min to 150~h; temperature
    $-20\,^{\circ}$C, $R = 84.4~\mu$m, $\alpha_c = 10^{-3}$; the numbered
    circles mark the frames in (a). Cite Demmenie et al.\ for the measured
    range and Appendix~B for the 0.29 prediction.
  \item \tG ``An isolated ice grain pair shown undergoing sintering shown'':
    ``Two equal ice grains sintering in saturated air, at $t = 0$ (left)
    and 150~h (right)''. ``filed'', ``the the''.
  \end{itemize}
\Pg{Figure 5, the gallery}{L(Section 4.1)}{One packing at each of five porosities, at four times, at $-20\,^{\circ}$C.}
  \begin{itemize}
  \item \tG ``Five microstructures of varying porosity $\varphi$ coarsening
    due to DSM'': ``One packing at each of five porosities, sintering at
    $-20\,^{\circ}$C''. Say which seed, and that each cell is 2~mm wide.
  \item \tA The vapor colorbar has words but no numbers: give the range, and
    that $\sigma$ is relative to flat ice and clipped, as in Figure~6.
  \item \tA The times (0, 10, 19, 30 d) are the dashed lines in
    Figure~7(a); one clause linking them.
  \end{itemize}
\Pg{Figure 6, one packing at five temperatures}{L(Section 4.2)}{(a) four frames, (b) conductivity against time, (c) against sintering age, (d) SSA against sintering age.}
  \begin{itemize}
  \item \tC (a): ``$|\sigma| \le 0.31$'': the colorbar says
    $0.34\times10^{-4}$. The factor $10^{-4}$ is missing from the caption.
  \item \tA Which packing (it appears to be the $\varphi = 0.325$ packing of
    Figure~5; say so) and that the frames are at $-20\,^{\circ}$C, at 0,
    10.2, 19.5 and 30.0 days.
  \item \tA Define $\theta = t/\tau_{\mathrm{sub}}$ and $\theta_r = 30$ (the
    dotted line). Say why the curves in (b) start at different values: each
    is normalized at $\theta_r$, which falls at a different time at each
    temperature.
  \item \tA (c), (d): the five curves overlie (within 0.3\%); that is the
    result, and the caption can say it. Axes start at $\theta = 1$
    (log scale).
  \item \tG ``preformed'' (performed); ``2.0mm'' (2.0~mm); ``The vertical
    axis is $k_{\mathrm{eff}}$, average value of $k_{xx}$ and $k_{yy}$, is
    normalized'': ``$k_{\mathrm{eff}} = (k_{xx}+k_{yy})/2$, normalized
    \ldots''. ``Shows the same curves plotted for'': ``(c) The same curves
    against''.
  \end{itemize}
\Pg{Figure 7, porosity}{L(Section 4.3)}{(a) conductivity against time, (b) against SSA, (c) its reference level, (d) its exponent.}
  \begin{itemize}
  \item \tC (a): ``The solid line corresponds to a single seed and the
    highlighted areas are the seed scatter'': the line is the mean of five
    packings and the band their minimum--maximum (as
    \texttt{compare\_keff.py} draws it).
    These are 2-D conductivities in W\,m$^{-1}$\,K$^{-1}$; say they are not
    those of 3-D snow (decided 09-23: no absolute value is claimed).
  \item \tC (b): only $\varphi \le 0.375$ is shown (three colors); the
    caption speaks of ``the porosities in panel (a)''. Say the higher
    porosities are left out and why (Section~4.3). $-40\,^{\circ}$C is
    absent because those runs end at $\theta = 43$; say so. The grey arrow
    is the direction of time.
  \item \tC (c), (d): one point per packing, temperatures pooled (they
    coincide at $\theta_r$), not ``each of the simulation runs''. The dashed
    line in (c) is a fit to $\varphi \le 0.375$, extended; give its form.
    The grey band marks the porosities left out of the fits.
  \item \tG ``Plots the the''; ``Shows'' at the start of (c), (d); ``15
    unique packings were used to model DSM at temperatures'': ``15 packings
    at each of $-30$, $-20$, $-10$ and $-5\,^{\circ}$C (60 curves)''.
  \end{itemize}
\end{enumerate}

\section*{Open Research; Acknowledgments}
\Miss Code and data availability; run inventory (\texttt{RUN\_TABLE.md}).
Acknowledgments.

\section*{Supporting Information}
\begin{itemize}
\item \textbf{Figures.} S1 domain-size convergence; S2 interface parameter;
  S3 time step and sampling; S4 temperature collapse for all packings; S5
  porosity trends, SSA by porosity, anisotropy in time.
\item \textbf{Tables, text.} Complete parameters, if moved; packing
  construction; the weak-form derivation, if trimmed; run table.
\end{itemize}

\section*{Appendix: the two calculations behind Sections 3 and 5.1}

\subsection*{A. Surface diffusion against vapor transport, by Molaro's own equations}

\textbf{Question.} Does surface diffusion supply the 59\% of neck growth our
model lacks at $-20\,^{\circ}$C?

\textbf{Equations} (Molaro et al., 2019, Table~3 and Eqs.~1--4, 6, 17--20;
$a$ grain radius, $x$ neck radius, $S$ their smoothing factor):
\begin{align*}
\dot V_1 &= S\,\frac{3\pi x\,\delta_sD_s\,\gamma_s\Omega\,(K_3 - K_2)}{d_2\,kT} && \text{surface diffusion}\\
\dot V_3 &= S\,2\pi x\rho\theta\,P_v\,\frac{\gamma_s\Omega}{kT}\left[\frac{\Omega}{2\pi\Delta_o kT}\right]^{1/2}(K_3 - K_m) && \text{vapor transport}\\
\dot x &= \frac{1}{2\pi x\theta\rho}\sum_i\dot V_i
\end{align*}
\[
K_3 = \frac{2}{a},\quad K_m = -\frac1\rho + \frac1x,\quad \rho = \frac{x^{2}}{2(a-x)},\quad \theta = \tan^{-1}\frac{a}{x+\rho},\quad K_1K_2 = K_m^{2},
\]
\[
\frac{K_2-K_1}{K_1} = \frac{3A\rho}{x} + \frac32 - \frac32\left[3\Big(\frac{A\rho}{x}\Big)^{2} + \frac{4A\rho}{x} + 1\right]^{1/2},\quad
d_2 = \frac{\rho}{1 + \left(\frac{4(K_2-K_1)}{3K_2}\right)^{1/2}},\quad A = \frac{\delta_bD_b}{\delta_sD_s} .
\]
The vapor term is Hertz--Knudsen with $\alpha = 1$ and no gas diffusion
(their Appendix~A).

\textbf{Parameters} (their Table~5): $\Omega = 1.181\times10^{-29}$~m$^{3}$;
$\Delta_o = 917$~kg\,m$^{-3}$; $\gamma_s = 0.06$~J\,m$^{-2}$; $\delta_s =
\delta_b = 9.04\times10^{-10}$~m; $D_{0s} = 1.4\times10^{-8}$~m$^{2}$\,s$^{-1}$
(literature range down to $1.35\times10^{-13}$), $Q_s = 23.2$~kJ\,mol$^{-1}$;
$D_{0b} = 8.4\times10^{-4}$~m$^{2}$\,s$^{-1}$, $Q_b = 4.9$~kJ\,mol$^{-1}$;
$P_v$ from Murphy \& Koop (2005), 103~Pa at $-20\,^{\circ}$C.

\textbf{At their experiment} ($a = 87~\mu$m, $-20\,^{\circ}$C; measured mean
$\dot x = 3.4\times10^{-9}$~m\,s$^{-1}$):
\begin{center}\small\begin{tabular}{rrrrrr}\toprule
$x$ [$\mu$m] & $\rho$ [$\mu$m] & $\dot V_1/\dot V_3$ & $\dot x$ from $\dot V_1$ [m/s] & $\dot x$ from $\dot V_3$ [m/s] & $\dot V_1/\dot V_3$, $A = 0$\\\midrule
16.4 & 1.9 & $6.7\times10^{-3}$ & $5\times10^{-11}$ & $7.8\times10^{-9}$ & $8\times10^{-7}$\\
20.0 & 3.0 & $3.1\times10^{-3}$ & $2\times10^{-11}$ & $4.9\times10^{-9}$ & $3\times10^{-7}$\\
25.0 & 5.0 & $1.3\times10^{-3}$ & $4\times10^{-12}$ & $2.9\times10^{-9}$ & $1\times10^{-7}$\\
32.4 & 9.6 & $4.1\times10^{-4}$ & $6\times10^{-13}$ & $1.5\times10^{-9}$ & $4\times10^{-8}$\\\bottomrule
\end{tabular}\end{center}
\begin{itemize}
\item Surface diffusion is 0.7\% falling to 0.04\% of their vapor term. It
  exceeds it only for $x$ below about 4~$\mu$m.
\item Their vapor term alone gives the measured rate.
\item The last column switches off the coupling to grain-boundary diffusion
  ($K_2 = K_m$). With their parameters $A = 3.5\times10^{8}$, and the
  coupling raises the driving curvature $K_2$ by a factor of about 4000.
\end{itemize}
\textbf{Check and its limit.} Their Figure~10 (25~$\mu$m, $T/T_m = 0.8$) was
recomputed. The vapor current agrees ($2\times10^{-21}$ rising to
$3\times10^{-20}$~m$^{3}$\,s$^{-1}$; theirs about $10^{-21}$ to
$2\times10^{-20}$). The surface current starts about 100 times above theirs,
because they solve the four coupled equations iteratively where this uses
the approximate Eq.~(20). The table therefore \emph{overstates} surface
diffusion. Script and output:
\texttt{studies/molaro\_2019/sa81\_mechanisms/}.

\textbf{What it implies.} The gap is between our vapor route ($\alpha_c =
0.1$, diffusion through air) and whatever supplies vapor in the experiment.
A one-equation neck model fitted to the same data wants $\alpha_c \approx
0.3$--1 (\texttt{studies/molaro\_2019/envelope\_ode/}).

\subsection*{B. The saturated pair against the Gibbs--Thomson law}

\textbf{Result (10-15).} From grains just touching, the run gives $a =
0.289$, on the prediction below (0.29). The tables that follow are for the
earlier \emph{pre-necked} run, kept as the record of why the initial
condition matters. \textbf{Question it answered:} was that run's 0.21
physics or resolution? Neither: it was the starting neck.

\textbf{Prediction.} For a neck of radius $x$ between grains of radius $R$,
with fillet radius $\rho = x^{2}/(2(R-x))$, fed from the grain surface
through attachment in series with gas diffusion over a length $c\rho$:
\[
v_n = \frac{d_0\left(\dfrac1\rho - \dfrac1x + \dfrac2R\right)}{\beta_{\mathrm{sub}} + c\,\rho\,k_D},\qquad
\frac{\mathrm dx}{\mathrm dt} = \frac{\pi}{2}\,v_n,\qquad k_D = \frac{\rho_{\mathrm{ice}}}{D_v\,\rho_{vs}^{I}} .
\]
Attachment alone gives $x^{3} \propto t$; diffusion alone $x^{5} \propto t$.
Why: with attachment limiting, $v_n \approx d_0/(\beta_{\mathrm{sub}}\rho)$
and $\rho \approx x^2/2R$, so $dx/dt \propto R/x^2$; with diffusion
limiting, the extra factor $1/(c\rho k_D)$ gives $dx/dt \propto R^2/x^4$.
The gas resistance grows with the neck; the attachment resistance does not.
This is Kingery's evaporation--condensation law ($n=3$) against Hobbs \&
Mason's diffusion through air ($n=5$; Blackford 2007, p.16).
$R = 84.4~\mu$m, $-20\,^{\circ}$C, $\alpha_c = 10^{-3}$; $c = 0.275$ from the
Molaro run, not refitted. Here gas diffusion is 0.4--2\% of the resistance.
Integrated from the run's first sample after relaxation, it gives a final
width of 60.6~$\mu$m (run: 63.2) and an exponent of 0.29 (run: 0.21).

\textbf{Run against prediction} ($\varepsilon = 0.12~\mu$m; one interface
band $= 9.19\,\varepsilon = 1.1~\mu$m; exponents from fits of
$C(t+t_0)^{a}$ starting at the time shown):
\begin{center}\small\begin{tabular}{rrrrrrr}\toprule
Start & $w$ [$\mu$m] & $w/\varepsilon$ & fillet [bands] & velocity, run / prediction & $a$, run & $a$, prediction\\\midrule
2 h & 35.1 & 293 & 2.1 & 2.54 & 0.215 & 0.292\\
6 h & 39.2 & 327 & 2.7 & 1.72 & 0.221 & 0.292\\
16 h & 44.9 & 374 & 3.7 & 1.36 & 0.227 & 0.290\\
31 h & 50.1 & 418 & 4.8 & 1.18 & 0.231 & 0.288\\
51 h & 55.0 & 459 & 6.0 & 1.08 & 0.234 & 0.287\\
71 h & 58.8 & 490 & 7.1 & 1.02 & 0.236 & 0.286\\\bottomrule
\end{tabular}\end{center}
The excess fits $1 + 11.4\,(\text{band}/\text{fillet})^{2.7}$ (rms 0.035).

\textbf{Sensitivity to the assumed contact time} (same data, $t_0$ fixed):
\begin{center}\small\begin{tabular}{lrrrrr}\toprule
$t_0$ & 4.8 h (free fit) & 8 h & 12 h & 16 h & 24 h\\\midrule
exponent & 0.212 & 0.234 & 0.258 & 0.281 & 0.325\\
rms misfit & 0.28\% & 0.78\% & 1.33\% & 1.73\% & 2.31\%\\\bottomrule
\end{tabular}\end{center}
The prediction takes 16~h to grow a neck to the run's first fitted width.

\textbf{What it shows and what it does not.} The rate converges on the
sharp-interface law, and the early excess accounts for most of the
difference in exponent. Two causes fit these numbers equally: the interface
being too wide for a small fillet, or the pre-made neck relaxing from its
imposed shape. The August run from tangent contact, at twice this interface
width, showed no excess and gave 0.283, which points to the pre-made neck.
The run from contact at this interface width (10-15) settles it:
$a = 0.289$, so the low exponent came from the pre-made neck.
From a pre-made neck the exponent does not reach the prediction within
100~h (0.236 at the end against 0.29); from contact it does at once. Still
to do: the velocity-over-prediction table for the tangent-start run. Scripts and output:
\texttt{studies/molaro\_2019/demmenie/} (\texttt{predict\_neck\_growth.py},
\texttt{interface\_width\_estimate.py}).

\subsection*{C. What the appendices of the manuscript should hold}
The draft now has two appendix headings with these titles. Suggested
content, each under a page:
\begin{itemize}
\item \emph{Appendix A, Surface diffusion against vapor transport:} the two
  rate equations and the curvature relations, the parameter list, the
  four-row table, and the two sentences on what was and was not reproduced
  of Molaro's Figure~10.
\item \emph{Appendix B, The saturated pair against the Gibbs--Thomson law:}
  the neck law, the run-against-prediction table, the contact-time table,
  led by the tangent-start run against the prediction (0.289 against
  0.29), then one paragraph and the contact-time table showing what a
  pre-made neck does to a fitted exponent.
\item The time--temperature scaling (two pages, in
  \texttt{temperature\_scaling/}) is a
  candidate third appendix, or supporting information, for Section~5.2.
\end{itemize}

\section*{Citations: what the .bib lacks}
Suggested above and present in the \texttt{.bib}: all of them. Mentioned in
this plan and \emph{not} in the \texttt{.bib}:
\begin{itemize}
\item Auriault, Boutin \& Geindreau (2009), the book whose naming of
  $\mathbf{t}$ the paper follows and which Calonne et al.\ (2014, p.6) cite
  for the cell problem (doi:10.1002/9780470612033; PDF in the library).
\item Ettrich et al.\ (2014), the tensor law applied to heat conduction;
  only if a second reference for the law is wanted.
\item A source for packing and percolation thresholds of discs, if
  Section~5.3 states them.
\end{itemize}
\Have Since 10-15 the \texttt{.bib} has Hobbs \& Mason (1964), Kingery (1960)
and Thomas et al.\ (1994); Kuczynski and Kingery \& Berg are no longer
needed, since Kingery (1960, Table~I) and Demmenie et al.\ (Table~1) give the
full set (F4). Blackford (2007) and Maeno \& Ebinuma (1983) are in it too.
Still to fix in the \texttt{.bib} (F11): Byrne et al.\ (2024) is the
preprint; \texttt{noauthor\_sintering\_nodate} duplicates
\texttt{Hobbs\_and\_Mason} and \texttt{noauthor\_strength\_nodate} duplicates
\texttt{choukroun\_strength\_2020}; old-style keys
\texttt{choukroun\_strength\_2020}, \texttt{nicoli\_tensorial\_2011},
\texttt{spitale\_association\_2007}, \texttt{postberg\_sodium\_2009},
\texttt{postberg\_macromolecular\_2018}, \texttt{kingery\_regelation\_1960},
\texttt{thomas\_crushing\_1994}, \texttt{Hobbs\_and\_Mason}.

\section*{Corrections that recur throughout}
Paragraph-specific corrections are under each paragraph. These apply
everywhere:
\begin{itemize}
\item \textbf{Notation.} Subscripts ``ice'' and ``air'' in place of $i$ and
  $a$ ($\rho_{\mathrm{ice}}$, $\phi_{\mathrm{ice}}$, $\phi_{\mathrm{air}}$);
  $\mathbf{K}_{\mathrm{eff}}$ for the tensor and $k_{\mathrm{eff}}$ for its
  scalar mean; $k_{\mathrm{ice}}$, $k_{\mathrm{air}}$ as in the tables (the
  text has $K_{\mathrm{ice}}$); $\rho_{\mathrm{air}}$, $P_{\mathrm{atm}}$ in
  Eq.~(2); $\phi$ against $\phi_i$ for the same field in Section~2.3.
\item \textbf{Temperature.} $T$ in $^\circ$C (user, 10-16), with
  $T + T_0$ ($T_0 = 273.15$~K) only in Eqs.~(3), (5), (7) and $d_0$;
  $T_0 = 273.15$~K only; the tables give the run temperature as $T$.
\item \textbf{$\varepsilon$} is the interface parameter, not the interface
  width: fixed in Section~2.1 (``controls the width''); one ``interface
  width'' remains (Section~2.3).
\item \textbf{``effect''} for ``affect'' (six places); ``preform'' for
  ``perform'' (four, incl.\ ``preformed'' in \P4); doubled words where \verb|\cref| already prints the
  word (``table table'', ``figure figure'', ``section section'').
\item \textbf{$\alpha_c$} written $1.002\times10^{-3}$, not ``1.002e-03''.
\item \textbf{Not cited in the text:} Table~2, Table~3, Figure~2.
\end{itemize}

\end{document}
